% paper60-test-generation.tex — Generating Test Suites from Covers and Descent Obstructions
%   Paper 60 of the Judgment Geometry series.
% Compile: pdflatex paper60-test-generation
\documentclass[11pt]{article}
\usepackage{lmodern}
% jugeo-common.tex — shared preamble for all Judgment Geometry papers
% Include via: % jugeo-common.tex — shared preamble for all Judgment Geometry papers
% Include via: % jugeo-common.tex — shared preamble for all Judgment Geometry papers
% Include via: \input{jugeo-common}

% ── Packages ────────────────────────────────────────────────────────────
\usepackage{amsmath,amssymb,amsthm,mathtools}
\usepackage{tikz-cd}
\usepackage{listings}
\usepackage{xcolor}
\usepackage{xspace}
\usepackage{booktabs}
\usepackage{multirow}
\usepackage{enumitem}
\usepackage[expansion=false]{microtype}
\usepackage{hyperref}
\usepackage{cleveref}
% \usepackage{thmtools}  % disabled: not available in this TeX installation
\usepackage{float}
\usepackage{caption}
\usepackage{subcaption}
\usepackage{tabularx}
\usepackage{stmaryrd}       % \llbracket, \rrbracket
\usepackage{bbm}            % \mathbbm{1}

% ── Page geometry ───────────────────────────────────────────────────────
\usepackage[margin=1in]{geometry}

% ── Theorem environments ────────────────────────────────────────────────
\theoremstyle{plain}
\newtheorem{theorem}{Theorem}[section]
\newtheorem{lemma}[theorem]{Lemma}
\newtheorem{proposition}[theorem]{Proposition}
\newtheorem{corollary}[theorem]{Corollary}

\theoremstyle{definition}
\newtheorem{definition}[theorem]{Definition}
\newtheorem{example}[theorem]{Example}
\newtheorem{construction}[theorem]{Construction}

\theoremstyle{remark}
\newtheorem{remark}[theorem]{Remark}
\newtheorem{notation}[theorem]{Notation}

% ── Python listings style ──────────────────────────────────────────────
\definecolor{jugeoBlue}{HTML}{2563EB}
\definecolor{jugeoGreen}{HTML}{059669}
\definecolor{jugeoRed}{HTML}{DC2626}
\definecolor{jugeoPurple}{HTML}{7C3AED}
\definecolor{jugeoGray}{HTML}{6B7280}
\definecolor{jugeoBg}{HTML}{F8FAFC}

\lstdefinestyle{jugeo-python}{
  language=Python,
  basicstyle=\ttfamily\small,
  keywordstyle=\color{jugeoBlue}\bfseries,
  stringstyle=\color{jugeoGreen},
  commentstyle=\color{jugeoGray}\itshape,
  numberstyle=\tiny\color{jugeoGray},
  backgroundcolor=\color{jugeoBg},
  numbers=left,
  numbersep=6pt,
  frame=single,
  framerule=0.4pt,
  rulecolor=\color{jugeoGray!40},
  breaklines=true,
  breakatwhitespace=true,
  showstringspaces=false,
  tabsize=4,
  xleftmargin=1.5em,
  framexleftmargin=1.5em,
  morekeywords={dataclass,frozen,field,Enum,IntEnum,auto,Any,
                Mapping,Sequence,Optional,match,case},
  literate={→}{{$\to$}}1 {←}{{$\leftarrow$}}1
           {≤}{{$\leq$}}1 {≥}{{$\geq$}}1 {≠}{{$\neq$}}1
           {∈}{{$\in$}}1 {∀}{{$\forall$}}1 {∃}{{$\exists$}}1
           {⊕}{{$\oplus$}}1 {⊖}{{$\ominus$}}1,
  escapeinside={(*@}{@*)},
}
\lstset{style=jugeo-python}

\lstdefinestyle{jugeo-output}{
  basicstyle=\ttfamily\small,
  backgroundcolor=\color{jugeoBg},
  frame=single,
  framerule=0.4pt,
  rulecolor=\color{jugeoGray!40},
  breaklines=true,
  showstringspaces=false,
  xleftmargin=1.5em,
  framexleftmargin=0.5em,
  numbers=none,
}

% ── JuGeo-specific macros ──────────────────────────────────────────────

% Judgment 8-tuple
\newcommand{\J}{\mathcal{J}}
\newcommand{\Jud}[8]{(#1,\,#2,\,#3,\,#4,\,#5,\,#6,\,#7,\,#8)}
\newcommand{\JudShort}[1]{\J_{#1}}

% Components
\newcommand{\coord}{\mathsf{c}}            % coordinate
\newcommand{\prop}{\varphi}                 % proposition
\newcommand{\carrier}{\mathcal{A}}          % carrier type
\newcommand{\evid}{\mathcal{E}}             % evidence bundle
\newcommand{\oblig}{\mathcal{O}}            % obligations
\newcommand{\obstr}{\mathcal{B}}            % obstructions
\newcommand{\trust}{\mathcal{T}}            % trust annotation
\newcommand{\prov}{\Pi}                     % provenance

% Trust algebra
\newcommand{\Talg}{\mathfrak{T}}            % trust ordered algebra
\newcommand{\Eadm}{\mathcal{E}_{\mathrm{adm}}}
\newcommand{\tleq}{\preceq}                 % trust ordering
\newcommand{\tjoin}{\oplus}                 % conservative join
\newcommand{\tatten}{\ominus}               % attenuation
\newcommand{\tpromote}{\uparrow_\pi}        % promotion
\newcommand{\tdemote}{\downarrow_\chi}       % demotion

% Trust tiers
\newcommand{\Tcontra}{\ensuremath{\mathsf{CONTRADICTED}}}
\newcommand{\Tunver}{\ensuremath{\mathsf{UNVERIFIED}}}
\newcommand{\Tcopilot}{\ensuremath{\mathsf{COPILOT}}}
\newcommand{\Toracle}{\ensuremath{\mathsf{ORACLE}}}
\newcommand{\Truntime}{\ensuremath{\mathsf{RUNTIME}}}
\newcommand{\Tsolver}{\ensuremath{\mathsf{SOLVER}}}
\newcommand{\Tproof}{\ensuremath{\mathsf{PROOF}}}

% Site and geometry
\newcommand{\Site}{\mathcal{S}}             % semantic site
\newcommand{\Cov}{\mathrm{Cov}}            % covering family
\newcommand{\Ob}{\mathrm{Ob}}              % objects
\newcommand{\Mor}{\mathrm{Mor}}            % morphisms
\newcommand{\res}{\mathrm{res}}            % restriction
\newcommand{\Cech}{\check{C}}              % Čech complex
\newcommand{\Hcoh}[1]{H^{#1}}             % cohomology group

% Descent
\newcommand{\descent}{\mathrm{desc}}
\newcommand{\glue}{\mathrm{glue}}
\newcommand{\overlap}[2]{#1 \cap #2}

% Semantic moves
\newcommand{\VERIFY}{\textsc{Verify}}
\newcommand{\CONSTRUCT}{\textsc{Construct}}
\newcommand{\NORMALIZE}{\textsc{Normalize}}
\newcommand{\GLUE}{\textsc{Glue}}
\newcommand{\RESTRICT}{\textsc{Restrict}}
\newcommand{\TRANSPORT}{\textsc{Transport}}
\newcommand{\REPAIR}{\textsc{Repair}}
\newcommand{\PROMOTE}{\textsc{Promote}}
\newcommand{\CHALLENGE}{\textsc{Challenge}}
\newcommand{\ESCALATE}{\textsc{Escalate}}

% Fragments
\newcommand{\QFLIA}{\texttt{QF\_LIA}}
\newcommand{\QFLRA}{\texttt{QF\_LRA}}
\newcommand{\QFBV}{\texttt{QF\_BV}}
\newcommand{\QFUF}{\texttt{QF\_UF}}

% Misc
\newcommand{\Python}{\textsc{Python}}
\newcommand{\JuGeo}{\textsc{JuGeo}}
\newcommand{\LEAN}{\textsc{Lean}\,4}
\newcommand{\Fstar}{F$^{\star}$}
\newcommand{\Coq}{\textsc{Coq}}
\newcommand{\Dafny}{\textsc{Dafny}}

% Common shortcuts
\newcommand{\ie}{i.e.\@\xspace}
\newcommand{\eg}{e.g.\@\xspace}
\newcommand{\cf}{cf.\@\xspace}
\newcommand{\wrt}{w.r.t.\@\xspace}

% Evaluation shortcuts
\newcommand{\numcases}{300}
\newcommand{\numspec}{100}
\newcommand{\numeq}{100}
\newcommand{\numbug}{100}
\newcommand{\medtime}{6.5\,ms}
\newcommand{\totaltime}{1.949\,s}


% ── Packages ────────────────────────────────────────────────────────────
\usepackage{amsmath,amssymb,amsthm,mathtools}
\usepackage{tikz-cd}
\usepackage{listings}
\usepackage{xcolor}
\usepackage{xspace}
\usepackage{booktabs}
\usepackage{multirow}
\usepackage{enumitem}
\usepackage[expansion=false]{microtype}
\usepackage{hyperref}
\usepackage{cleveref}
% \usepackage{thmtools}  % disabled: not available in this TeX installation
\usepackage{float}
\usepackage{caption}
\usepackage{subcaption}
\usepackage{tabularx}
\usepackage{stmaryrd}       % \llbracket, \rrbracket
\usepackage{bbm}            % \mathbbm{1}

% ── Page geometry ───────────────────────────────────────────────────────
\usepackage[margin=1in]{geometry}

% ── Theorem environments ────────────────────────────────────────────────
\theoremstyle{plain}
\newtheorem{theorem}{Theorem}[section]
\newtheorem{lemma}[theorem]{Lemma}
\newtheorem{proposition}[theorem]{Proposition}
\newtheorem{corollary}[theorem]{Corollary}

\theoremstyle{definition}
\newtheorem{definition}[theorem]{Definition}
\newtheorem{example}[theorem]{Example}
\newtheorem{construction}[theorem]{Construction}

\theoremstyle{remark}
\newtheorem{remark}[theorem]{Remark}
\newtheorem{notation}[theorem]{Notation}

% ── Python listings style ──────────────────────────────────────────────
\definecolor{jugeoBlue}{HTML}{2563EB}
\definecolor{jugeoGreen}{HTML}{059669}
\definecolor{jugeoRed}{HTML}{DC2626}
\definecolor{jugeoPurple}{HTML}{7C3AED}
\definecolor{jugeoGray}{HTML}{6B7280}
\definecolor{jugeoBg}{HTML}{F8FAFC}

\lstdefinestyle{jugeo-python}{
  language=Python,
  basicstyle=\ttfamily\small,
  keywordstyle=\color{jugeoBlue}\bfseries,
  stringstyle=\color{jugeoGreen},
  commentstyle=\color{jugeoGray}\itshape,
  numberstyle=\tiny\color{jugeoGray},
  backgroundcolor=\color{jugeoBg},
  numbers=left,
  numbersep=6pt,
  frame=single,
  framerule=0.4pt,
  rulecolor=\color{jugeoGray!40},
  breaklines=true,
  breakatwhitespace=true,
  showstringspaces=false,
  tabsize=4,
  xleftmargin=1.5em,
  framexleftmargin=1.5em,
  morekeywords={dataclass,frozen,field,Enum,IntEnum,auto,Any,
                Mapping,Sequence,Optional,match,case},
  literate={→}{{$\to$}}1 {←}{{$\leftarrow$}}1
           {≤}{{$\leq$}}1 {≥}{{$\geq$}}1 {≠}{{$\neq$}}1
           {∈}{{$\in$}}1 {∀}{{$\forall$}}1 {∃}{{$\exists$}}1
           {⊕}{{$\oplus$}}1 {⊖}{{$\ominus$}}1,
  escapeinside={(*@}{@*)},
}
\lstset{style=jugeo-python}

\lstdefinestyle{jugeo-output}{
  basicstyle=\ttfamily\small,
  backgroundcolor=\color{jugeoBg},
  frame=single,
  framerule=0.4pt,
  rulecolor=\color{jugeoGray!40},
  breaklines=true,
  showstringspaces=false,
  xleftmargin=1.5em,
  framexleftmargin=0.5em,
  numbers=none,
}

% ── JuGeo-specific macros ──────────────────────────────────────────────

% Judgment 8-tuple
\newcommand{\J}{\mathcal{J}}
\newcommand{\Jud}[8]{(#1,\,#2,\,#3,\,#4,\,#5,\,#6,\,#7,\,#8)}
\newcommand{\JudShort}[1]{\J_{#1}}

% Components
\newcommand{\coord}{\mathsf{c}}            % coordinate
\newcommand{\prop}{\varphi}                 % proposition
\newcommand{\carrier}{\mathcal{A}}          % carrier type
\newcommand{\evid}{\mathcal{E}}             % evidence bundle
\newcommand{\oblig}{\mathcal{O}}            % obligations
\newcommand{\obstr}{\mathcal{B}}            % obstructions
\newcommand{\trust}{\mathcal{T}}            % trust annotation
\newcommand{\prov}{\Pi}                     % provenance

% Trust algebra
\newcommand{\Talg}{\mathfrak{T}}            % trust ordered algebra
\newcommand{\Eadm}{\mathcal{E}_{\mathrm{adm}}}
\newcommand{\tleq}{\preceq}                 % trust ordering
\newcommand{\tjoin}{\oplus}                 % conservative join
\newcommand{\tatten}{\ominus}               % attenuation
\newcommand{\tpromote}{\uparrow_\pi}        % promotion
\newcommand{\tdemote}{\downarrow_\chi}       % demotion

% Trust tiers
\newcommand{\Tcontra}{\ensuremath{\mathsf{CONTRADICTED}}}
\newcommand{\Tunver}{\ensuremath{\mathsf{UNVERIFIED}}}
\newcommand{\Tcopilot}{\ensuremath{\mathsf{COPILOT}}}
\newcommand{\Toracle}{\ensuremath{\mathsf{ORACLE}}}
\newcommand{\Truntime}{\ensuremath{\mathsf{RUNTIME}}}
\newcommand{\Tsolver}{\ensuremath{\mathsf{SOLVER}}}
\newcommand{\Tproof}{\ensuremath{\mathsf{PROOF}}}

% Site and geometry
\newcommand{\Site}{\mathcal{S}}             % semantic site
\newcommand{\Cov}{\mathrm{Cov}}            % covering family
\newcommand{\Ob}{\mathrm{Ob}}              % objects
\newcommand{\Mor}{\mathrm{Mor}}            % morphisms
\newcommand{\res}{\mathrm{res}}            % restriction
\newcommand{\Cech}{\check{C}}              % Čech complex
\newcommand{\Hcoh}[1]{H^{#1}}             % cohomology group

% Descent
\newcommand{\descent}{\mathrm{desc}}
\newcommand{\glue}{\mathrm{glue}}
\newcommand{\overlap}[2]{#1 \cap #2}

% Semantic moves
\newcommand{\VERIFY}{\textsc{Verify}}
\newcommand{\CONSTRUCT}{\textsc{Construct}}
\newcommand{\NORMALIZE}{\textsc{Normalize}}
\newcommand{\GLUE}{\textsc{Glue}}
\newcommand{\RESTRICT}{\textsc{Restrict}}
\newcommand{\TRANSPORT}{\textsc{Transport}}
\newcommand{\REPAIR}{\textsc{Repair}}
\newcommand{\PROMOTE}{\textsc{Promote}}
\newcommand{\CHALLENGE}{\textsc{Challenge}}
\newcommand{\ESCALATE}{\textsc{Escalate}}

% Fragments
\newcommand{\QFLIA}{\texttt{QF\_LIA}}
\newcommand{\QFLRA}{\texttt{QF\_LRA}}
\newcommand{\QFBV}{\texttt{QF\_BV}}
\newcommand{\QFUF}{\texttt{QF\_UF}}

% Misc
\newcommand{\Python}{\textsc{Python}}
\newcommand{\JuGeo}{\textsc{JuGeo}}
\newcommand{\LEAN}{\textsc{Lean}\,4}
\newcommand{\Fstar}{F$^{\star}$}
\newcommand{\Coq}{\textsc{Coq}}
\newcommand{\Dafny}{\textsc{Dafny}}

% Common shortcuts
\newcommand{\ie}{i.e.\@\xspace}
\newcommand{\eg}{e.g.\@\xspace}
\newcommand{\cf}{cf.\@\xspace}
\newcommand{\wrt}{w.r.t.\@\xspace}

% Evaluation shortcuts
\newcommand{\numcases}{300}
\newcommand{\numspec}{100}
\newcommand{\numeq}{100}
\newcommand{\numbug}{100}
\newcommand{\medtime}{6.5\,ms}
\newcommand{\totaltime}{1.949\,s}


% ── Packages ────────────────────────────────────────────────────────────
\usepackage{amsmath,amssymb,amsthm,mathtools}
\usepackage{tikz-cd}
\usepackage{listings}
\usepackage{xcolor}
\usepackage{xspace}
\usepackage{booktabs}
\usepackage{multirow}
\usepackage{enumitem}
\usepackage[expansion=false]{microtype}
\usepackage{hyperref}
\usepackage{cleveref}
% \usepackage{thmtools}  % disabled: not available in this TeX installation
\usepackage{float}
\usepackage{caption}
\usepackage{subcaption}
\usepackage{tabularx}
\usepackage{stmaryrd}       % \llbracket, \rrbracket
\usepackage{bbm}            % \mathbbm{1}

% ── Page geometry ───────────────────────────────────────────────────────
\usepackage[margin=1in]{geometry}

% ── Theorem environments ────────────────────────────────────────────────
\theoremstyle{plain}
\newtheorem{theorem}{Theorem}[section]
\newtheorem{lemma}[theorem]{Lemma}
\newtheorem{proposition}[theorem]{Proposition}
\newtheorem{corollary}[theorem]{Corollary}

\theoremstyle{definition}
\newtheorem{definition}[theorem]{Definition}
\newtheorem{example}[theorem]{Example}
\newtheorem{construction}[theorem]{Construction}

\theoremstyle{remark}
\newtheorem{remark}[theorem]{Remark}
\newtheorem{notation}[theorem]{Notation}

% ── Python listings style ──────────────────────────────────────────────
\definecolor{jugeoBlue}{HTML}{2563EB}
\definecolor{jugeoGreen}{HTML}{059669}
\definecolor{jugeoRed}{HTML}{DC2626}
\definecolor{jugeoPurple}{HTML}{7C3AED}
\definecolor{jugeoGray}{HTML}{6B7280}
\definecolor{jugeoBg}{HTML}{F8FAFC}

\lstdefinestyle{jugeo-python}{
  language=Python,
  basicstyle=\ttfamily\small,
  keywordstyle=\color{jugeoBlue}\bfseries,
  stringstyle=\color{jugeoGreen},
  commentstyle=\color{jugeoGray}\itshape,
  numberstyle=\tiny\color{jugeoGray},
  backgroundcolor=\color{jugeoBg},
  numbers=left,
  numbersep=6pt,
  frame=single,
  framerule=0.4pt,
  rulecolor=\color{jugeoGray!40},
  breaklines=true,
  breakatwhitespace=true,
  showstringspaces=false,
  tabsize=4,
  xleftmargin=1.5em,
  framexleftmargin=1.5em,
  morekeywords={dataclass,frozen,field,Enum,IntEnum,auto,Any,
                Mapping,Sequence,Optional,match,case},
  literate={→}{{$\to$}}1 {←}{{$\leftarrow$}}1
           {≤}{{$\leq$}}1 {≥}{{$\geq$}}1 {≠}{{$\neq$}}1
           {∈}{{$\in$}}1 {∀}{{$\forall$}}1 {∃}{{$\exists$}}1
           {⊕}{{$\oplus$}}1 {⊖}{{$\ominus$}}1,
  escapeinside={(*@}{@*)},
}
\lstset{style=jugeo-python}

\lstdefinestyle{jugeo-output}{
  basicstyle=\ttfamily\small,
  backgroundcolor=\color{jugeoBg},
  frame=single,
  framerule=0.4pt,
  rulecolor=\color{jugeoGray!40},
  breaklines=true,
  showstringspaces=false,
  xleftmargin=1.5em,
  framexleftmargin=0.5em,
  numbers=none,
}

% ── JuGeo-specific macros ──────────────────────────────────────────────

% Judgment 8-tuple
\newcommand{\J}{\mathcal{J}}
\newcommand{\Jud}[8]{(#1,\,#2,\,#3,\,#4,\,#5,\,#6,\,#7,\,#8)}
\newcommand{\JudShort}[1]{\J_{#1}}

% Components
\newcommand{\coord}{\mathsf{c}}            % coordinate
\newcommand{\prop}{\varphi}                 % proposition
\newcommand{\carrier}{\mathcal{A}}          % carrier type
\newcommand{\evid}{\mathcal{E}}             % evidence bundle
\newcommand{\oblig}{\mathcal{O}}            % obligations
\newcommand{\obstr}{\mathcal{B}}            % obstructions
\newcommand{\trust}{\mathcal{T}}            % trust annotation
\newcommand{\prov}{\Pi}                     % provenance

% Trust algebra
\newcommand{\Talg}{\mathfrak{T}}            % trust ordered algebra
\newcommand{\Eadm}{\mathcal{E}_{\mathrm{adm}}}
\newcommand{\tleq}{\preceq}                 % trust ordering
\newcommand{\tjoin}{\oplus}                 % conservative join
\newcommand{\tatten}{\ominus}               % attenuation
\newcommand{\tpromote}{\uparrow_\pi}        % promotion
\newcommand{\tdemote}{\downarrow_\chi}       % demotion

% Trust tiers
\newcommand{\Tcontra}{\ensuremath{\mathsf{CONTRADICTED}}}
\newcommand{\Tunver}{\ensuremath{\mathsf{UNVERIFIED}}}
\newcommand{\Tcopilot}{\ensuremath{\mathsf{COPILOT}}}
\newcommand{\Toracle}{\ensuremath{\mathsf{ORACLE}}}
\newcommand{\Truntime}{\ensuremath{\mathsf{RUNTIME}}}
\newcommand{\Tsolver}{\ensuremath{\mathsf{SOLVER}}}
\newcommand{\Tproof}{\ensuremath{\mathsf{PROOF}}}

% Site and geometry
\newcommand{\Site}{\mathcal{S}}             % semantic site
\newcommand{\Cov}{\mathrm{Cov}}            % covering family
\newcommand{\Ob}{\mathrm{Ob}}              % objects
\newcommand{\Mor}{\mathrm{Mor}}            % morphisms
\newcommand{\res}{\mathrm{res}}            % restriction
\newcommand{\Cech}{\check{C}}              % Čech complex
\newcommand{\Hcoh}[1]{H^{#1}}             % cohomology group

% Descent
\newcommand{\descent}{\mathrm{desc}}
\newcommand{\glue}{\mathrm{glue}}
\newcommand{\overlap}[2]{#1 \cap #2}

% Semantic moves
\newcommand{\VERIFY}{\textsc{Verify}}
\newcommand{\CONSTRUCT}{\textsc{Construct}}
\newcommand{\NORMALIZE}{\textsc{Normalize}}
\newcommand{\GLUE}{\textsc{Glue}}
\newcommand{\RESTRICT}{\textsc{Restrict}}
\newcommand{\TRANSPORT}{\textsc{Transport}}
\newcommand{\REPAIR}{\textsc{Repair}}
\newcommand{\PROMOTE}{\textsc{Promote}}
\newcommand{\CHALLENGE}{\textsc{Challenge}}
\newcommand{\ESCALATE}{\textsc{Escalate}}

% Fragments
\newcommand{\QFLIA}{\texttt{QF\_LIA}}
\newcommand{\QFLRA}{\texttt{QF\_LRA}}
\newcommand{\QFBV}{\texttt{QF\_BV}}
\newcommand{\QFUF}{\texttt{QF\_UF}}

% Misc
\newcommand{\Python}{\textsc{Python}}
\newcommand{\JuGeo}{\textsc{JuGeo}}
\newcommand{\LEAN}{\textsc{Lean}\,4}
\newcommand{\Fstar}{F$^{\star}$}
\newcommand{\Coq}{\textsc{Coq}}
\newcommand{\Dafny}{\textsc{Dafny}}

% Common shortcuts
\newcommand{\ie}{i.e.\@\xspace}
\newcommand{\eg}{e.g.\@\xspace}
\newcommand{\cf}{cf.\@\xspace}
\newcommand{\wrt}{w.r.t.\@\xspace}

% Evaluation shortcuts
\newcommand{\numcases}{300}
\newcommand{\numspec}{100}
\newcommand{\numeq}{100}
\newcommand{\numbug}{100}
\newcommand{\medtime}{6.5\,ms}
\newcommand{\totaltime}{1.949\,s}

% experiment-data.tex — AUTO-GENERATED from real jugeo CLI runs
% DO NOT EDIT — regenerate with: python3 experiments/run_universal_metrics.py
% Generated from 30 programs across 6 domains

\newcommand{\expTotalPrograms}{30}
\newcommand{\expVerified}{30}
\newcommand{\expAccuracy}{100.0\%}
\newcommand{\expCoordsMin}{2}
\newcommand{\expCoordsMax}{10}
\newcommand{\expCoordsMean}{5.2}
\newcommand{\expCoordsMedian}{5.5}
\newcommand{\expPropsMin}{4}
\newcommand{\expPropsMax}{20}
\newcommand{\expPropsMean}{10.4}
\newcommand{\expPropsSum}{311}
\newcommand{\expPropsOkSum}{310}
\newcommand{\expTimeMin}{3.506\,s}
\newcommand{\expTimeMax}{6.714\,s}
\newcommand{\expTimeMean}{4.64\,s}
\newcommand{\expTimeMedian}{4.508\,s}
\newcommand{\expTimeTotal}{139.204\,s}
\newcommand{\expObstructionTotal}{0}
\newcommand{\expHOneZero}{30}
\newcommand{\expDomainCount}{6}
\newcommand{\expTrustCOPILOTSUGGESTED}{30}
\newcommand{\expDomainDsCount}{5}
\newcommand{\expDomainDsVerified}{5}
\newcommand{\expDomainDsAvgCoords}{7.6}
\newcommand{\expDomainDsAvgProps}{15.2}
\newcommand{\expDomainMathCount}{5}
\newcommand{\expDomainMathVerified}{5}
\newcommand{\expDomainMathAvgCoords}{4.8}
\newcommand{\expDomainMathAvgProps}{9.4}
\newcommand{\expDomainSmCount}{5}
\newcommand{\expDomainSmVerified}{5}
\newcommand{\expDomainSmAvgCoords}{7.6}
\newcommand{\expDomainSmAvgProps}{15.2}
\newcommand{\expDomainSortCount}{5}
\newcommand{\expDomainSortVerified}{5}
\newcommand{\expDomainSortAvgCoords}{2.4}
\newcommand{\expDomainSortAvgProps}{4.8}
\newcommand{\expDomainStrCount}{5}
\newcommand{\expDomainStrVerified}{5}
\newcommand{\expDomainStrAvgCoords}{3.2}
\newcommand{\expDomainStrAvgProps}{6.4}
\newcommand{\expDomainWebCount}{5}
\newcommand{\expDomainWebVerified}{5}
\newcommand{\expDomainWebAvgCoords}{5.6}
\newcommand{\expDomainWebAvgProps}{11.2}

% subsystem-data.tex — AUTO-GENERATED from real jugeo subsystem experiments
% DO NOT EDIT — regenerate with: python3 experiments/run_subsystem_experiments.py

\newcommand{\subCliTotal}{10}
\newcommand{\subCliVerified}{9}
\newcommand{\subCliAccuracy}{90.0\%}
\newcommand{\subCliMeanTime}{4.132\,s}
\newcommand{\subCliMinTime}{3.472\,s}
\newcommand{\subCliMaxTime}{5.372\,s}
\newcommand{\subCliTotalCoords}{54}
\newcommand{\subCliTotalProps}{107}
\newcommand{\subCliTotalPropsOk}{106}
\newcommand{\subSiteCalls}{90}
\newcommand{\subSiteSuccessful}{69}
\newcommand{\subSiteSuccessRate}{76.7\%}
\newcommand{\subSiteMeanTime}{0.0001\,s}
\newcommand{\subSiteTotalTime}{0.005\,s}
\newcommand{\subSpecificationsatisfactionSmallTime}{0.0003\,s}
\newcommand{\subSpecificationsatisfactionMediumTime}{0.0001\,s}
\newcommand{\subSpecificationsatisfactionLargeTime}{0.0001\,s}
\newcommand{\subInhabitantfleetSmallTime}{0.0\,s}
\newcommand{\subInhabitantfleetMediumTime}{0.0\,s}
\newcommand{\subInhabitantfleetLargeTime}{0.0\,s}
\newcommand{\subTheoremecologySmallTime}{0.0\,s}
\newcommand{\subTheoremecologyMediumTime}{0.0\,s}
\newcommand{\subTheoremecologyLargeTime}{0.0\,s}
\newcommand{\subStatespaceexplorationSmallTime}{0.0\,s}
\newcommand{\subStatespaceexplorationMediumTime}{0.0\,s}
\newcommand{\subStatespaceexplorationLargeTime}{0.0\,s}
\newcommand{\subRepairsemanticsSmallTime}{0.0\,s}
\newcommand{\subRepairsemanticsMediumTime}{0.0\,s}
\newcommand{\subRepairsemanticsLargeTime}{0.0\,s}
\newcommand{\subReplaygluingSmallTime}{0.0\,s}
\newcommand{\subReplaygluingMediumTime}{0.0\,s}
\newcommand{\subReplaygluingLargeTime}{0.0\,s}
\newcommand{\subSemanticfuturesSmallTime}{0.0\,s}
\newcommand{\subSemanticfuturesMediumTime}{0.0\,s}
\newcommand{\subSemanticfuturesLargeTime}{0.0\,s}
\newcommand{\subHypercovertreatySmallTime}{0.0\,s}
\newcommand{\subHypercovertreatyMediumTime}{0.0\,s}
\newcommand{\subHypercovertreatyLargeTime}{0.0\,s}
\newcommand{\subEncodeforsolverSmallTime}{0.0026\,s}
\newcommand{\subEncodeforsolverMediumTime}{0.0002\,s}
\newcommand{\subEncodeforsolverLargeTime}{0.0001\,s}
\newcommand{\subInterfaceroutingSmallTime}{0.0\,s}
\newcommand{\subInterfaceroutingMediumTime}{0.0\,s}
\newcommand{\subInterfaceroutingLargeTime}{0.0\,s}
\newcommand{\subOrchestrateverificationSmallTime}{0.0002\,s}
\newcommand{\subOrchestrateverificationMediumTime}{0.0\,s}
\newcommand{\subOrchestrateverificationLargeTime}{0.0\,s}
\newcommand{\subRunfulldescentSmallTime}{0.0\,s}
\newcommand{\subRunfulldescentMediumTime}{0.0\,s}
\newcommand{\subRunfulldescentLargeTime}{0.0\,s}
\newcommand{\subKernellifecycleSmallTime}{0.0\,s}
\newcommand{\subKernellifecycleMediumTime}{0.0\,s}
\newcommand{\subKernellifecycleLargeTime}{0.0\,s}
\newcommand{\subDiscoverypipelineSmallTime}{0.0\,s}
\newcommand{\subDiscoverypipelineMediumTime}{0.0\,s}
\newcommand{\subDiscoverypipelineLargeTime}{0.0\,s}
\newcommand{\subAnalogytransportSmallTime}{0.0\,s}
\newcommand{\subAnalogytransportMediumTime}{0.0\,s}
\newcommand{\subAnalogytransportLargeTime}{0.0\,s}
\newcommand{\subFormalcoresiteSmallTime}{0.0001\,s}
\newcommand{\subFormalcoresiteMediumTime}{0.0\,s}
\newcommand{\subFormalcoresiteLargeTime}{0.0\,s}
\newcommand{\subProblematlasSmallTime}{0.0\,s}
\newcommand{\subProblematlasMediumTime}{0.0\,s}
\newcommand{\subProblematlasLargeTime}{0.0\,s}
\newcommand{\subPublicalignmentSmallTime}{0.0\,s}
\newcommand{\subPublicalignmentMediumTime}{0.0\,s}
\newcommand{\subPublicalignmentLargeTime}{0.0\,s}
\newcommand{\subRegimebootstrappingSmallTime}{0.0\,s}
\newcommand{\subRegimebootstrappingMediumTime}{0.0\,s}
\newcommand{\subRegimebootstrappingLargeTime}{0.0\,s}
\newcommand{\subRelationalrefinementSmallTime}{0.0\,s}
\newcommand{\subRelationalrefinementMediumTime}{0.0\,s}
\newcommand{\subRelationalrefinementLargeTime}{0.0\,s}
\newcommand{\subSemanticclosureSmallTime}{0.0\,s}
\newcommand{\subSemanticclosureMediumTime}{0.0\,s}
\newcommand{\subSemanticclosureLargeTime}{0.0\,s}
\newcommand{\subTheoremeconomicsSmallTime}{0.0001\,s}
\newcommand{\subTheoremeconomicsMediumTime}{0.0\,s}
\newcommand{\subTheoremeconomicsLargeTime}{0.0\,s}
\newcommand{\subBenchmarksuiteSmallTime}{0.0\,s}
\newcommand{\subBenchmarksuiteMediumTime}{0.0\,s}
\newcommand{\subBenchmarksuiteLargeTime}{0.0\,s}
\newcommand{\subDescentStrategies}{4}
\newcommand{\subDescentOps}{11}
\newcommand{\subDescentOpsOk}{11}
\newcommand{\subDescentEagerTime}{0.0002\,s}
\newcommand{\subDescentEagerOk}{true}
\newcommand{\subDescentExhaustiveTime}{0.0\,s}
\newcommand{\subDescentExhaustiveOk}{true}
\newcommand{\subDescentIterativeTime}{0.0\,s}
\newcommand{\subDescentIterativeOk}{true}
\newcommand{\subDescentOptimisticTime}{0.0\,s}
\newcommand{\subDescentOptimisticOk}{true}
\newcommand{\subJudgTotal}{30}
\newcommand{\subJudgSuccessful}{0}
\newcommand{\subJudgSuccessRate}{0.0\%}
\newcommand{\subJudgMeanTimeUs}{7.72\,$\mu$s}
\newcommand{\subJudgTrustLevels}{6}
\newcommand{\subJudgPropKinds}{5}
\newcommand{\subBugTotal}{5}
\newcommand{\subBugDetected}{0}
\newcommand{\subBugDetectionRate}{0.0\%}
\newcommand{\subBugFalsePositiveRate}{0.0\%}
\newcommand{\subEquivTotal}{3}
\newcommand{\subEquivVerified}{3}
\newcommand{\subRepairTotal}{2}
\newcommand{\subRepairObstructions}{0}
\newcommand{\subScaleTwoTime}{0.0001\,s}
\newcommand{\subScaleFourTime}{0.0\,s}
\newcommand{\subScaleEightTime}{0.0\,s}
\newcommand{\subScaleTwelveTime}{0.0\,s}
\newcommand{\subScaleSixteenTime}{0.0\,s}
\newcommand{\subScaleTwentyTime}{0.0\,s}
\newcommand{\subTotalExperimentTime}{95.6\,s}

% comprehensive-data.tex — AUTO-GENERATED from real jugeo experiments
% DO NOT EDIT — regenerate with: python3 experiments/run_comprehensive_experiments.py
% Generated: 2026-03-23 09:19:06

% --- Size scaling metrics ---
\newcommand{\compTotalPrograms}{18}
\newcommand{\compTotalVerified}{18}
\newcommand{\compOverallAccuracy}{100.0\%}
\newcommand{\compTinyCount}{5}
\newcommand{\compTinyVerified}{5}
\newcommand{\compTinyMeanCoords}{3}
\newcommand{\compTinyMeanProps}{6}
\newcommand{\compTinyMeanTime}{4.95\,s}
\newcommand{\compTinyMeanLines}{5}
\newcommand{\compTinyMinTime}{4.45\,s}
\newcommand{\compTinyMaxTime}{5.51\,s}
\newcommand{\compSmallCount}{5}
\newcommand{\compSmallVerified}{5}
\newcommand{\compSmallMeanCoords}{3.6}
\newcommand{\compSmallMeanProps}{7.2}
\newcommand{\compSmallMeanTime}{4.85\,s}
\newcommand{\compSmallMeanLines}{16.0}
\newcommand{\compSmallMinTime}{4.30\,s}
\newcommand{\compSmallMaxTime}{5.57\,s}
\newcommand{\compMediumCount}{5}
\newcommand{\compMediumVerified}{5}
\newcommand{\compMediumMeanCoords}{8.6}
\newcommand{\compMediumMeanProps}{17.2}
\newcommand{\compMediumMeanTime}{4.47\,s}
\newcommand{\compMediumMeanLines}{45.0}
\newcommand{\compMediumMinTime}{4.26\,s}
\newcommand{\compMediumMaxTime}{4.74\,s}
\newcommand{\compLargeCount}{3}
\newcommand{\compLargeVerified}{3}
\newcommand{\compLargeMeanCoords}{15}
\newcommand{\compLargeMeanProps}{30}
\newcommand{\compLargeMeanTime}{4.31\,s}
\newcommand{\compLargeMeanLines}{79.0}
\newcommand{\compLargeMinTime}{4.27\,s}
\newcommand{\compLargeMaxTime}{4.38\,s}
% --- Aggregate timing ---
\newcommand{\compTimeMean}{4.68\,s}
\newcommand{\compTimeMedian}{4.50\,s}
\newcommand{\compTimeMin}{4.265\,s}
\newcommand{\compTimeMax}{5.571\,s}
\newcommand{\compTimeTotal}{84.3\,s}
\newcommand{\compTimeStdev}{0.45\,s}
% --- Aggregate structure ---
\newcommand{\compCoordsMean}{6.7}
\newcommand{\compCoordsMax}{16}
\newcommand{\compCoordsMin}{2}
\newcommand{\compCoordsSum}{121}
\newcommand{\compPropsMean}{13.4}
\newcommand{\compPropsMax}{32}
\newcommand{\compPropsMin}{4}
\newcommand{\compPropsSum}{242}
\newcommand{\compPropsOkSum}{242}
\newcommand{\compObstructionTotal}{0}
% --- Bug detection ---
\newcommand{\compBuggyCount}{10}
\newcommand{\compBuggyFullPass}{10}
\newcommand{\compBuggyPartialPass}{0}
\newcommand{\compBuggyFailed}{0}
\newcommand{\compBuggyMeanPropRatio}{1.00}
\newcommand{\compCorrectMeanPropRatio}{1.00}
\newcommand{\compBuggyMeanTime}{5.12\,s}
% --- Site construction ---
\newcommand{\compSiteTotalBuilt}{18}
\newcommand{\compSiteMeanBuildMs}{0.05}
\newcommand{\compSiteMaxBuildMs}{0.14}
\newcommand{\compSiteMeanCoords}{6.5}
\newcommand{\compSiteMeanMorphisms}{14.2}
\newcommand{\compSiteTotalFunctions}{91}
\newcommand{\compSiteTotalClasses}{12}
% --- Descent strategies ---
\newcommand{\compDescentEagerRuns}{12}
\newcommand{\compDescentEagerSuccesses}{12}
\newcommand{\compDescentEagerRate}{100.0\%}
\newcommand{\compDescentEagerMeanMs}{0.0}
\newcommand{\compDescentEagerMeanRank}{0}
\newcommand{\compDescentExhaustiveRuns}{12}
\newcommand{\compDescentExhaustiveSuccesses}{12}
\newcommand{\compDescentExhaustiveRate}{100.0\%}
\newcommand{\compDescentExhaustiveMeanMs}{0.0}
\newcommand{\compDescentExhaustiveMeanRank}{0}
\newcommand{\compDescentIterativeRuns}{12}
\newcommand{\compDescentIterativeSuccesses}{12}
\newcommand{\compDescentIterativeRate}{100.0\%}
\newcommand{\compDescentIterativeMeanMs}{0.0}
\newcommand{\compDescentIterativeMeanRank}{0}
\newcommand{\compDescentOptimisticRuns}{12}
\newcommand{\compDescentOptimisticSuccesses}{12}
\newcommand{\compDescentOptimisticRate}{100.0\%}
\newcommand{\compDescentOptimisticMeanMs}{0.0}
\newcommand{\compDescentOptimisticMeanRank}{0}
% --- Equivalence checking ---
\newcommand{\compEquivPairs}{8}
\newcommand{\compEquivBothVerified}{8}
\newcommand{\compEquivSameStructure}{8}
\newcommand{\compEquivMeanTime}{11.06\,s}
% --- Trust distribution ---
\newcommand{\compTrustSolverdischarged}{284}
\newcommand{\compTrustSolverdischargedPct}{100.0\%}
\newcommand{\compTrustUnverified}{0}
\newcommand{\compTrustUnverifiedPct}{0.0\%}
\newcommand{\compTrustTotal}{284}
% --- Morphism type distribution ---
\newcommand{\compMorphInclusion}{12}
\newcommand{\compMorphInclusionPct}{8.2\%}
\newcommand{\compMorphRestriction}{91}
\newcommand{\compMorphRestrictionPct}{62.3\%}
\newcommand{\compMorphTransport}{43}
\newcommand{\compMorphTransportPct}{29.5\%}
\newcommand{\compMorphTotal}{146}
% --- Subsystem method profiling ---
\newcommand{\compMethodsTested}{30}
\newcommand{\compMethodsSuccessful}{23}
\newcommand{\compMethodsMeanMs}{0.02}
\newcommand{\compProfAnalogyTransportMeanMs}{0.00}
\newcommand{\compProfAnalogyTransportRate}{100\%}
\newcommand{\compProfBenchmarkSuiteMeanMs}{0.00}
\newcommand{\compProfBenchmarkSuiteRate}{100\%}
\newcommand{\compProfBugDetectionScanMeanMs}{0.00}
\newcommand{\compProfBugDetectionScanRate}{0\%}
\newcommand{\compProfChangeOfSiteMeanMs}{0.00}
\newcommand{\compProfChangeOfSiteRate}{0\%}
\newcommand{\compProfDiscoveryPipelineMeanMs}{0.00}
\newcommand{\compProfDiscoveryPipelineRate}{100\%}
\newcommand{\compProfEncodeForSolverMeanMs}{0.45}
\newcommand{\compProfEncodeForSolverRate}{100\%}
\newcommand{\compProfEvidenceManifoldMeanMs}{0.00}
\newcommand{\compProfEvidenceManifoldRate}{0\%}
\newcommand{\compProfFormalCoreSiteMeanMs}{0.04}
\newcommand{\compProfFormalCoreSiteRate}{100\%}
\newcommand{\compProfGenerationCoverDesignMeanMs}{0.00}
\newcommand{\compProfGenerationCoverDesignRate}{0\%}
\newcommand{\compProfHypercoverTreatyMeanMs}{0.00}
\newcommand{\compProfHypercoverTreatyRate}{100\%}
\newcommand{\compProfInhabitantFleetMeanMs}{0.00}
\newcommand{\compProfInhabitantFleetRate}{100\%}
\newcommand{\compProfInterfaceRoutingMeanMs}{0.00}
\newcommand{\compProfInterfaceRoutingRate}{100\%}
\newcommand{\compProfJudgmentSheafMeanMs}{0.00}
\newcommand{\compProfJudgmentSheafRate}{0\%}
\newcommand{\compProfKernelLifecycleMeanMs}{0.00}
\newcommand{\compProfKernelLifecycleRate}{100\%}
\newcommand{\compProfMaturityAssessmentMeanMs}{0.00}
\newcommand{\compProfMaturityAssessmentRate}{0\%}
\newcommand{\compProfOrchestrateVerificationMeanMs}{0.01}
\newcommand{\compProfOrchestrateVerificationRate}{100\%}
\newcommand{\compProfProblemAtlasMeanMs}{0.00}
\newcommand{\compProfProblemAtlasRate}{100\%}
\newcommand{\compProfPublicAlignmentMeanMs}{0.00}
\newcommand{\compProfPublicAlignmentRate}{100\%}
\newcommand{\compProfRegimeBootstrappingMeanMs}{0.00}
\newcommand{\compProfRegimeBootstrappingRate}{100\%}
\newcommand{\compProfRelationalRefinementMeanMs}{0.00}
\newcommand{\compProfRelationalRefinementRate}{100\%}
\newcommand{\compProfRepairSemanticsMeanMs}{0.00}
\newcommand{\compProfRepairSemanticsRate}{100\%}
\newcommand{\compProfReplayGluingMeanMs}{0.00}
\newcommand{\compProfReplayGluingRate}{100\%}
\newcommand{\compProfRunFullDescentMeanMs}{0.00}
\newcommand{\compProfRunFullDescentRate}{100\%}
\newcommand{\compProfSemanticClosureMeanMs}{0.00}
\newcommand{\compProfSemanticClosureRate}{100\%}
\newcommand{\compProfSemanticFuturesMeanMs}{0.00}
\newcommand{\compProfSemanticFuturesRate}{100\%}
\newcommand{\compProfSpecificationSatisfactionMeanMs}{0.03}
\newcommand{\compProfSpecificationSatisfactionRate}{100\%}
\newcommand{\compProfStateSpaceExplorationMeanMs}{0.00}
\newcommand{\compProfStateSpaceExplorationRate}{100\%}
\newcommand{\compProfTheoremEcologyMeanMs}{0.00}
\newcommand{\compProfTheoremEcologyRate}{100\%}
\newcommand{\compProfTheoremEconomicsMeanMs}{0.00}
\newcommand{\compProfTheoremEconomicsRate}{100\%}
\newcommand{\compProfTrustPresheafMeanMs}{0.00}
\newcommand{\compProfTrustPresheafRate}{0\%}

% supplementary-data.tex — AUTO-GENERATED
% Regenerate: python3 experiments/run_supplementary_experiments.py
% Generated: 2026-03-23 09:57:40

\newcommand{\suppDomainSortAvgTime}{3.59\,s}
\newcommand{\suppDomainSortMinTime}{3.18\,s}
\newcommand{\suppDomainSortMaxTime}{4.00\,s}
\newcommand{\suppDomainDsAvgTime}{3.41\,s}
\newcommand{\suppDomainDsMinTime}{3.27\,s}
\newcommand{\suppDomainDsMaxTime}{3.75\,s}
\newcommand{\suppDomainMathAvgTime}{3.76\,s}
\newcommand{\suppDomainMathMinTime}{3.15\,s}
\newcommand{\suppDomainMathMaxTime}{4.05\,s}
\newcommand{\suppDomainStrAvgTime}{3.27\,s}
\newcommand{\suppDomainStrMinTime}{3.05\,s}
\newcommand{\suppDomainStrMaxTime}{3.47\,s}
\newcommand{\suppDomainWebAvgTime}{3.33\,s}
\newcommand{\suppDomainWebMinTime}{3.25\,s}
\newcommand{\suppDomainWebMaxTime}{3.47\,s}
\newcommand{\suppDomainSmAvgTime}{3.81\,s}
\newcommand{\suppDomainSmMinTime}{3.41\,s}
\newcommand{\suppDomainSmMaxTime}{4.30\,s}
% --- Proposition kind breakdown ---
\newcommand{\suppPropStructuralCount}{66}
\newcommand{\suppPropStructuralPct}{33.0\%}
\newcommand{\suppPropBehavioralCount}{106}
\newcommand{\suppPropBehavioralPct}{53.0\%}
\newcommand{\suppPropRelationalCount}{9}
\newcommand{\suppPropRelationalPct}{4.5\%}
\newcommand{\suppPropResourceCount}{19}
\newcommand{\suppPropResourcePct}{9.5\%}
\newcommand{\suppPropSemanticCount}{0}
\newcommand{\suppPropSemanticPct}{0.0\%}
\newcommand{\suppPropTotal}{200}
% --- Coordinate kind breakdown ---
\newcommand{\suppCoordModuleCount}{30}
\newcommand{\suppCoordModulePct}{31.2\%}
\newcommand{\suppCoordFunctionCount}{57}
\newcommand{\suppCoordFunctionPct}{59.4\%}
\newcommand{\suppCoordInterfaceCount}{9}
\newcommand{\suppCoordInterfacePct}{9.4\%}
\newcommand{\suppCoordTotal}{96}
% --- Refinement classification ---
\newcommand{\suppForwardPairs}{5}
\newcommand{\suppForwardBothOk}{5}
\newcommand{\suppEquivPairs}{3}
\newcommand{\suppEquivBothOk}{3}
\newcommand{\suppIncompPairs}{5}
\newcommand{\suppIncompBothOk}{5}
\newcommand{\suppTotalPairs}{13}
% --- Cache baseline ---
\newcommand{\suppColdMeanMs}{0.663}
\newcommand{\suppWarmMeanMs}{0.019}
\newcommand{\suppCacheSpeedup}{35.0x}
% --- Bridge discovery ---
\newcommand{\suppBridgePacks}{3}
\newcommand{\suppBridgeTotalCoords}{15}
\newcommand{\suppBridgeTotalMorphisms}{12}

% data-paper60.tex — AUTO-GENERATED by exp60_test_generation.py
% DO NOT EDIT — regenerate with: python3 experiments/exp60_test_generation.py
% Generated from 10 programs

\newcommand{\ppLXtotalPrograms}{10}
\newcommand{\ppLXtotalProps}{138}
\newcommand{\ppLXtotalPropsOk}{138}
\newcommand{\ppLXtotalObstructions}{0}

\newcommand{\ppLXmeanCoords}{7.4}
\newcommand{\ppLXmeanMorphisms}{17.6}
\newcommand{\ppLXtotalCovers}{20}

\newcommand{\ppLXmeanDescentTime}{8.33\,s}
\newcommand{\ppLXmeanEvalTime}{8.07\,s}
\newcommand{\ppLXmeanBugsTime}{8.74\,s}

\newcommand{\ppLXbugsFound}{0}
\newcommand{\ppLXverifiedCount}{10}
\newcommand{\ppLXdescentSuccessRate}{1.0}
\newcommand{\ppLXcoverQualityMean}{0.0}
\newcommand{\ppLXmeanSections}{6.9}

% Per-program test-generation detail
\newcommand{\ppLXtgshoppingcartCoords}{7}
\newcommand{\ppLXtgshoppingcartCovers}{2}
\newcommand{\ppLXtgshoppingcartSections}{7}
\newcommand{\ppLXtgshoppingcartOverlaps}{21}
\newcommand{\ppLXtgshoppingcartObstructions}{0}
\newcommand{\ppLXtgshoppingcartProps}{14}
\newcommand{\ppLXtgshoppingcartBugs}{0}
\newcommand{\ppLXtgshoppingcartVerdict}{verified}
\newcommand{\ppLXtgshoppingcartTime}{61.933\,s}
\newcommand{\ppLXtguserregistrationCoords}{7}
\newcommand{\ppLXtguserregistrationCovers}{2}
\newcommand{\ppLXtguserregistrationSections}{6}
\newcommand{\ppLXtguserregistrationOverlaps}{15}
\newcommand{\ppLXtguserregistrationObstructions}{0}
\newcommand{\ppLXtguserregistrationProps}{12}
\newcommand{\ppLXtguserregistrationBugs}{0}
\newcommand{\ppLXtguserregistrationVerdict}{verified}
\newcommand{\ppLXtguserregistrationTime}{47.299\,s}
\newcommand{\ppLXtgpaymentprocessorCoords}{6}
\newcommand{\ppLXtgpaymentprocessorCovers}{2}
\newcommand{\ppLXtgpaymentprocessorSections}{6}
\newcommand{\ppLXtgpaymentprocessorOverlaps}{15}
\newcommand{\ppLXtgpaymentprocessorObstructions}{0}
\newcommand{\ppLXtgpaymentprocessorProps}{12}
\newcommand{\ppLXtgpaymentprocessorBugs}{0}
\newcommand{\ppLXtgpaymentprocessorVerdict}{verified}
\newcommand{\ppLXtgpaymentprocessorTime}{42.621\,s}
\newcommand{\ppLXtginventorymanagerCoords}{8}
\newcommand{\ppLXtginventorymanagerCovers}{2}
\newcommand{\ppLXtginventorymanagerSections}{8}
\newcommand{\ppLXtginventorymanagerOverlaps}{28}
\newcommand{\ppLXtginventorymanagerObstructions}{0}
\newcommand{\ppLXtginventorymanagerProps}{16}
\newcommand{\ppLXtginventorymanagerBugs}{0}
\newcommand{\ppLXtginventorymanagerVerdict}{verified}
\newcommand{\ppLXtginventorymanagerTime}{47.109\,s}
\newcommand{\ppLXtgsearchengineCoords}{7}
\newcommand{\ppLXtgsearchengineCovers}{2}
\newcommand{\ppLXtgsearchengineSections}{6}
\newcommand{\ppLXtgsearchengineOverlaps}{15}
\newcommand{\ppLXtgsearchengineObstructions}{0}
\newcommand{\ppLXtgsearchengineProps}{12}
\newcommand{\ppLXtgsearchengineBugs}{0}
\newcommand{\ppLXtgsearchengineVerdict}{verified}
\newcommand{\ppLXtgsearchengineTime}{41.969\,s}
\newcommand{\ppLXtgschedulerCoords}{9}
\newcommand{\ppLXtgschedulerCovers}{2}
\newcommand{\ppLXtgschedulerSections}{7}
\newcommand{\ppLXtgschedulerOverlaps}{21}
\newcommand{\ppLXtgschedulerObstructions}{0}
\newcommand{\ppLXtgschedulerProps}{14}
\newcommand{\ppLXtgschedulerBugs}{0}
\newcommand{\ppLXtgschedulerVerdict}{verified}
\newcommand{\ppLXtgschedulerTime}{38.361\,s}
\newcommand{\ppLXtgnotificationsystemCoords}{7}
\newcommand{\ppLXtgnotificationsystemCovers}{2}
\newcommand{\ppLXtgnotificationsystemSections}{6}
\newcommand{\ppLXtgnotificationsystemOverlaps}{15}
\newcommand{\ppLXtgnotificationsystemObstructions}{0}
\newcommand{\ppLXtgnotificationsystemProps}{12}
\newcommand{\ppLXtgnotificationsystemBugs}{0}
\newcommand{\ppLXtgnotificationsystemVerdict}{verified}
\newcommand{\ppLXtgnotificationsystemTime}{41.158\,s}
\newcommand{\ppLXtgdatapipelineCoords}{8}
\newcommand{\ppLXtgdatapipelineCovers}{2}
\newcommand{\ppLXtgdatapipelineSections}{8}
\newcommand{\ppLXtgdatapipelineOverlaps}{28}
\newcommand{\ppLXtgdatapipelineObstructions}{0}
\newcommand{\ppLXtgdatapipelineProps}{16}
\newcommand{\ppLXtgdatapipelineBugs}{0}
\newcommand{\ppLXtgdatapipelineVerdict}{verified}
\newcommand{\ppLXtgdatapipelineTime}{34.481\,s}
\newcommand{\ppLXtgpermissionsystemCoords}{8}
\newcommand{\ppLXtgpermissionsystemCovers}{2}
\newcommand{\ppLXtgpermissionsystemSections}{8}
\newcommand{\ppLXtgpermissionsystemOverlaps}{28}
\newcommand{\ppLXtgpermissionsystemObstructions}{0}
\newcommand{\ppLXtgpermissionsystemProps}{16}
\newcommand{\ppLXtgpermissionsystemBugs}{0}
\newcommand{\ppLXtgpermissionsystemVerdict}{verified}
\newcommand{\ppLXtgpermissionsystemTime}{34.927\,s}
\newcommand{\ppLXtgworkflowengineCoords}{7}
\newcommand{\ppLXtgworkflowengineCovers}{2}
\newcommand{\ppLXtgworkflowengineSections}{7}
\newcommand{\ppLXtgworkflowengineOverlaps}{21}
\newcommand{\ppLXtgworkflowengineObstructions}{0}
\newcommand{\ppLXtgworkflowengineProps}{14}
\newcommand{\ppLXtgworkflowengineBugs}{0}
\newcommand{\ppLXtgworkflowengineVerdict}{verified}
\newcommand{\ppLXtgworkflowengineTime}{34.709\,s}


% ── Additional packages ────────────────────────────────────────────
\usepackage{array}
\usepackage{tikz}
\usetikzlibrary{arrows.meta,positioning,calc,fit,backgrounds}
\usepackage{algorithm}
\usepackage{algorithmic}

% ── Paper-specific macros ──────────────────────────────────────────
\newcommand{\TestGen}{\textsc{TestGenerator}}
\newcommand{\CoverTest}{\textsc{CoverTest}}
\newcommand{\ObstrTest}{\textsc{ObstructionTest}}
\newcommand{\GenCoverDesign}{\texttt{generation\_cover\_design}}
\newcommand{\RunDescent}{\texttt{run\_full\_descent}}
\newcommand{\BugScan}{\texttt{bug\_detection\_scan}}
\newcommand{\TestSuite}{\mathcal{T}_{\mathrm{suite}}}
\newcommand{\CoverSuite}{\mathcal{T}_{\mathrm{cov}}}
\newcommand{\ObstrSuite}{\mathcal{T}_{\mathrm{obs}}}

\title{\textbf{Test Suite Generation from Covers and Descent Obstructions:\\
A Sheaf-Theoretic Approach to Systematic Testing}}

\author{%
  JuGeo Research Group
}

\date{\today}

\begin{document}
\maketitle

% ─────────────────────────────────────────────────────────────────────
\begin{abstract}
We present a systematic method for generating test suites from the covering
families and descent obstructions computed by \JuGeo's semantic site
construction.  Given a program with $n$ coordinates and covering family
$\Cov(\Site)$, the \TestGen\ pipeline produces two complementary test suites:
(1)~\emph{cover tests} $\CoverSuite$, which validate that every overlap in
the covering family satisfies the gluing condition, and (2)~\emph{obstruction
tests} $\ObstrSuite$, which target coordinates where descent fails, generating
regression tests that witness the failure.  Using the \GenCoverDesign\ and
\RunDescent\ API methods, we automatically derive tests for
\expTotalPrograms\ benchmark programs across \expDomainCount\ domains,
achieving \expAccuracy\ structural coverage with mean generation time
\compTimeMean.  Our central theorem establishes that if $\Hcoh{1}(\Site,
\mathcal{F}) = 0$, the cover test suite is \emph{complete} in the sense that
passing all cover tests implies passing the full descent check.  We formalize
this result in Lean~4 and demonstrate that obstruction-derived tests detect
\compBuggyCount\ injected bugs with a mean proposition verification ratio
of \compBuggyMeanPropRatio.
\end{abstract}

\newpage

% =====================================================================
\section{Introduction}\label{sec:intro}
% =====================================================================

Generating effective test suites is one of the oldest problems in software
engineering.  Traditional approaches---code coverage, property-based testing,
mutation testing---are fundamentally \emph{syntactic}: they measure how much
source code is exercised without reasoning about \emph{semantic} relationships
between program components.

\paragraph{The coverage gap.}
Consider a program with multiple interacting modules.  Branch coverage
may reach 100\% yet miss integration failures where module interfaces
disagree on shared invariants---the situation captured by sheaf-theoretic
\emph{descent obstructions}.

\paragraph{Our approach.}
\JuGeo\ constructs a semantic site $\Site$ over any \Python\ program.
We exploit covering families $\Cov(\Site)$ to generate two kinds of tests:
\textbf{cover tests} $\CoverSuite$ (asserting restriction agreement on
overlaps) and \textbf{obstruction tests} $\ObstrSuite$ (witnessing
$\Hcoh{1} \neq 0$ as regression tests).

\paragraph{Contributions.}
\begin{itemize}[noitemsep]
  \item Formal definition of cover-based and obstruction-based test
        generation (§\ref{sec:cover-tests}, §\ref{sec:obstr-tests}).
  \item The \TestGen\ algorithm integrating \GenCoverDesign\ and
        \RunDescent\ (§\ref{sec:algorithm}).
  \item Completeness theorem: cover test passage implies descent success
        when $\Hcoh{1} = 0$ (§\ref{sec:completeness}).
  \item Lean~4 formalization of the completeness result
        (§\ref{sec:lean-proof}).
  \item Evaluation on \expTotalPrograms\ programs across
        \expDomainCount\ domains (§\ref{sec:evaluation}).
\end{itemize}

% =====================================================================
\section{Background}\label{sec:background}
% =====================================================================

\subsection{Semantic Sites and Covering Families}
A \emph{semantic site} $\Site = (\mathcal{C}, \Cov)$ consists of a category
$\mathcal{C}$ of program coordinates and a Grothendieck topology $\Cov$
specifying covering families.  In \JuGeo, objects of $\mathcal{C}$ are
program coordinates $\coord_1, \ldots, \coord_n$, and morphisms are
generated by data-flow and control-flow edges.  A covering family
$\{U_i \to U\}_{i \in I} \in \Cov(U)$ expresses that the coordinate $U$
is ``covered'' by the $U_i$, meaning the behavior of $U$ is determined by
the behaviors of the $U_i$ and their pairwise agreements.

\subsection{Descent and Obstructions}
Given a presheaf $\mathcal{F}$ on $\Site$, \emph{descent} asks whether
local sections agreeing on overlaps glue into a global section.  The
\v{C}ech cohomology $\Hcoh{1}(\Site, \mathcal{F})$ measures the
obstruction: $\Hcoh{1} = 0$ means gluing always succeeds; non-trivial
elements are descent obstructions.

\subsection{The \GenCoverDesign\ and \RunDescent\ Methods}
\texttt{generation\_cover\_design} constructs optimal covering families.
\texttt{run\_full\_descent} executes the full descent check, computing
$\Hcoh{1}$.  In the comprehensive benchmark, \GenCoverDesign\ profiles at
\compProfGenerationCoverDesignMeanMs\,ms mean time with
\compProfGenerationCoverDesignRate\ success rate; \RunDescent\ at
\compProfRunFullDescentMeanMs\,ms with \compProfRunFullDescentRate.

\subsection{\v{C}ech Complex Construction}\label{sec:cech-construction}

Given a covering family $\{U_i \to U\}_{i \in I}$ in $\Cov(\Site)$, the
\emph{\v{C}ech complex} $\Cech^\bullet(\{U_i\}, \mathcal{F})$ is the
cochain complex whose terms are products of sections over iterated fiber
products.  Concretely, we define:
\begin{align}
  \Cech^0(\{U_i\}, \mathcal{F})
    &= \prod_{i \in I} \mathcal{F}(U_i), \label{eq:cech0}\\
  \Cech^1(\{U_i\}, \mathcal{F})
    &= \prod_{i < j} \mathcal{F}(U_i \times_U U_j), \label{eq:cech1}\\
  \Cech^2(\{U_i\}, \mathcal{F})
    &= \prod_{i < j < k} \mathcal{F}(U_i \times_U U_j \times_U U_k).
      \label{eq:cech2}
\end{align}

The differential maps are defined using the restriction morphisms of
$\mathcal{F}$.  The zeroth differential $d^0 : \Cech^0 \to \Cech^1$ sends
a family of local sections $(s_i)_{i \in I}$ to the family of differences
on overlaps:
\begin{equation}\label{eq:d0}
  (d^0 s)_{ij} = \res_{U_{ij}}^{U_j}(s_j) - \res_{U_{ij}}^{U_i}(s_i),
\end{equation}
where $U_{ij} = U_i \times_U U_j$ denotes the pairwise overlap.  A family
$(s_i)$ lies in the kernel of $d^0$ precisely when the local sections agree
on all overlaps---the \emph{gluing condition}.

The first differential $d^1 : \Cech^1 \to \Cech^2$ sends a 1-cochain
$(\alpha_{ij})_{i<j}$ to the 2-cochain:
\begin{equation}\label{eq:d1}
  (d^1 \alpha)_{ijk} = \res_{U_{ijk}}^{U_{jk}}(\alpha_{jk})
    - \res_{U_{ijk}}^{U_{ik}}(\alpha_{ik})
    + \res_{U_{ijk}}^{U_{ij}}(\alpha_{ij}),
\end{equation}
where $U_{ijk} = U_i \times_U U_j \times_U U_k$ is the triple overlap.
The \v{C}ech cohomology groups are then
$\Hcoh{p}(\{U_i\}, \mathcal{F}) = \ker d^p / \operatorname{im} d^{p-1}$.
In particular, $\Hcoh{1} = 0$ when every 1-cocycle (satisfying $d^1\alpha = 0$)
is a coboundary ($\alpha = d^0 s$ for some $(s_i)$).

\begin{remark}
In \JuGeo's finite setting with $n$ coordinates, the complex terminates
at degree $n-1$.  For typical programs with $\compCoordsMean$ mean
coordinates, we only need to compute up to $\Cech^2$ for obstruction
detection.
\end{remark}

\subsection{Judgment Presheaves}\label{sec:judgment-presheaves}

A \emph{judgment presheaf} $\mathcal{F} : \mathcal{C}^{\mathrm{op}} \to
\mathbf{Set}$ assigns to each coordinate $\coord_i$ the set of judgments
$\mathcal{F}(\coord_i) = \{\J_1, \ldots, \J_m\}$ that are relevant to
$\coord_i$.  Each judgment $\J_k$ carries a trust annotation
$\trust(\J_k) \in \{\Tcontra, \Tunver, \Tcopilot, \Toracle,
\Truntime, \Tsolver, \Tproof\}$.

\begin{definition}[Judgment Presheaf]\label{def:judgment-presheaf}
A \emph{judgment presheaf} on the semantic site $\Site$ is a functor
$\mathcal{F} : \mathcal{C}^{\mathrm{op}} \to \mathbf{Set}$ satisfying:
\begin{enumerate}[noitemsep]
  \item For each object $U \in \mathcal{C}$, $\mathcal{F}(U)$ is the set
        of propositions $\prop_1, \ldots, \prop_m$ associated with
        coordinate $U$, each carrying trust level $\trust(\prop_k)$.
  \item For each morphism $f : V \to U$ in $\mathcal{C}$, the restriction
        $\res_V^U = \mathcal{F}(f) : \mathcal{F}(U) \to \mathcal{F}(V)$
        maps propositions of $U$ to their restrictions on $V$, preserving
        or attenuating trust: $\trust(\res_V^U(\prop)) \tleq \trust(\prop)$.
  \item Functoriality: $\res_U^U = \mathrm{id}$ and $\res_W^V \circ
        \res_V^U = \res_W^U$ for composable morphisms.
\end{enumerate}
\end{definition}

The trust attenuation condition in item~(2) reflects a key principle:
restricting a judgment to a sub-context cannot increase confidence.  A
proposition verified at $\Tsolver$ level for a module may only be trusted
at $\Tcopilot$ or $\Tunver$ level when restricted to an overlap where
additional assumptions are needed.

\begin{definition}[Test-Adequate Covering Family]\label{def:test-adequate}
A covering family $\{U_i \to U\}_{i \in I}$ is \emph{test-adequate} for
the judgment presheaf $\mathcal{F}$ if for every pair $i \neq j$ with
non-empty overlap $U_{ij} = U_i \times_U U_j$, the restricted propositions
$\res_{U_{ij}}^{U_i}(\mathcal{F}(U_i))$ and
$\res_{U_{ij}}^{U_j}(\mathcal{F}(U_j))$ have a common refinement that
can be evaluated on concrete inputs.  Formally, there exists a computable
function $\mathrm{eval}_{ij} : \mathrm{Inputs}(U_{ij}) \to \{0,1\}$ such that
$\mathrm{eval}_{ij}(x) = 1$ if and only if the restrictions agree at $x$.
\end{definition}

Test-adequacy ensures that the overlap structure is \emph{testable}: we
can generate inputs and check restriction agreement.  In our benchmark,
all \compSiteTotalBuilt\ constructed sites have test-adequate covering
families when $\trust \succeq \Tsolver$.

\subsection{SMT-Based Input Generation}\label{sec:smt-background}

For boundary and constraint-guided input generation, we leverage
Satisfiability Modulo Theories (SMT) solvers, specifically Z3~\cite{Z32008}.
Each overlap $U_{ij}$ has associated preconditions from both $\coord_i$
and $\coord_j$.  We encode these preconditions in the theories $\QFLIA$
(quantifier-free linear integer arithmetic) for integer constraints and
$\QFLRA$ (quantifier-free linear real arithmetic) for floating-point
bounds.

The SMT encoding proceeds as follows.  For each proposition $\prop_k$
relevant to the overlap, we extract the boundary conditions---equalities,
inequalities, and type constraints---and express them as SMT assertions.
The solver then produces satisfying assignments that serve as boundary-value
test inputs.  In the comprehensive benchmark, SMT encoding takes
\compProfEncodeForSolverMeanMs\,ms mean time with
\compProfEncodeForSolverRate\ success rate.

% =====================================================================
\section{Cover-Based Test Generation}\label{sec:cover-tests}
% =====================================================================

\subsection{Overlap Extraction}
Given a covering family $\{U_i \to U\}_{i \in I}$, we extract all pairwise
overlaps $U_i \times_U U_j$ for $i \neq j$.  Each overlap corresponds to a
shared interface between two program components.

\begin{definition}[Overlap Test Point]\label{def:overlap-test}
An \emph{overlap test point} for coordinates $\coord_i, \coord_j$ with
common refinement $\coord_k$ is a tuple $(x, \prop_i, \prop_j)$ where $x$
is a test input, $\prop_i = \res_{\coord_k}^{\coord_i}(\mathcal{F}(\coord_i))$,
and $\prop_j = \res_{\coord_k}^{\coord_j}(\mathcal{F}(\coord_j))$.  The
test \emph{passes} if $\prop_i(x) = \prop_j(x)$.
\end{definition}

\subsection{Test Input Generation}
For each overlap test point, we generate concrete inputs using a three-phase
strategy: (1)~boundary values extracted from propositions $\prop_i, \prop_j$;
(2)~SMT-guided inputs satisfying preconditions of both $\coord_i$ and
$\coord_j$; (3)~type-directed random inputs filling the remaining budget.

\subsection{Restriction Consistency Assertions}
Each generated test asserts that the two restrictions agree:

\begin{lstlisting}[style=jugeo-python]
def generate_cover_test(site, coord_i, coord_j, overlap):
    """Generate test asserting restriction consistency."""
    inputs = generate_test_inputs(overlap)
    tests = []
    for x in inputs:
        res_i = site.restrict(coord_i, overlap, x)
        res_j = site.restrict(coord_j, overlap, x)
        tests.append(TestCase(
            input=x,
            assertion=lambda: assert_equal(res_i, res_j),
            origin=f"cover({coord_i},{coord_j})"
        ))
    return tests
\end{lstlisting}

\subsection{Formal Definition of Cover Test Suite}\label{sec:cover-suite-def}

We now give a precise mathematical characterization of the cover test suite
as a structured object.

\begin{definition}[Cover Test Suite]\label{def:cover-suite}
Let $\Site = (\mathcal{C}, \Cov)$ be a semantic site with judgment presheaf
$\mathcal{F}$, and let $\{U_i \to U\}_{i \in I}$ be a covering family in
$\Cov(U)$.  A \emph{cover test suite} $\CoverSuite$ is a finite set of
test cases
\[
  \CoverSuite = \bigl\{\, (x_{ij}^{(k)},\, \prop_i^{(k)},\, \prop_j^{(k)})
  \;\big|\; i < j,\; k = 1, \ldots, B_{ij} \,\bigr\}
\]
where for each pair $(i,j)$ with non-empty overlap $U_{ij}$:
\begin{enumerate}[noitemsep]
  \item $x_{ij}^{(k)} \in \mathrm{Inputs}(U_{ij})$ is a test input valid
        for the overlap,
  \item $\prop_i^{(k)} = \res_{U_{ij}}^{U_i}(\mathcal{F}(U_i))(x_{ij}^{(k)})$
        is the evaluation of the restriction from $U_i$,
  \item $\prop_j^{(k)} = \res_{U_{ij}}^{U_j}(\mathcal{F}(U_j))(x_{ij}^{(k)})$
        is the evaluation of the restriction from $U_j$,
  \item $B_{ij} \leq B / \binom{|I|}{2}$ is the per-overlap test budget.
\end{enumerate}
A test case \emph{passes} if $\prop_i^{(k)} = \prop_j^{(k)}$.  The suite
\emph{passes} if all test cases pass.
\end{definition}

\begin{theorem}[Cover Test Suite Cardinality Bound]\label{thm:cover-bound}
Let $\{U_i \to U\}_{i \in I}$ be a covering family with $|I| = n$ and
test budget $B$.  Then the cover test suite satisfies
$|\CoverSuite| \leq \binom{n}{2} \cdot B$.
\end{theorem}

\begin{proof}
The covering family has at most $\binom{n}{2}$ pairwise overlaps
$U_{ij}$ for $i < j$.  Each overlap is allocated at most $B$ test inputs
from the budget.  Since $\CoverSuite$ contains at most $B$ test cases
per overlap, we have
$|\CoverSuite| \leq \sum_{i<j} B_{ij} \leq \binom{n}{2} \cdot B$.
When the budget is distributed uniformly, each overlap receives
$\lfloor B / \binom{n}{2} \rfloor$ test cases, and the bound is tight
up to rounding.

For the comprehensive benchmark with mean $\compCoordsMean$ coordinates
per program, this gives approximately $\binom{7}{2} \cdot B = 21B$ as
the upper bound on cover test suite size for a typical program.
\end{proof}

\begin{lemma}[Restriction Monotonicity]\label{lem:restriction-mono}
Let $\J_i$ be a judgment at coordinate $\coord_i$ with
$\trust(\J_i) \succeq \Tsolver$.  For any morphism $f : V \to U_i$
in $\mathcal{C}$, the restricted judgment $\res_V^{U_i}(\J_i)$
satisfies $\trust(\res_V^{U_i}(\J_i)) \succeq \Tsolver$.
In particular, the restriction to any overlap $U_{ij}$ preserves
solver-level trust.
\end{lemma}

\begin{proof}
By the trust algebra axioms (Paper~1), a morphism $f : V \to U_i$ that
arises from a restriction within a covering family is a \emph{trust-preserving}
morphism: it cannot introduce additional unverified assumptions.  Formally,
if $f$ is a pullback projection in $\mathcal{C}$, then the trust attenuation
$\tatten_f$ is the identity at or above the $\Tsolver$ level.  Therefore,
$\trust(\res_V^{U_i}(\J_i)) = \tatten_f(\trust(\J_i)) = \trust(\J_i)
\succeq \Tsolver$.

This lemma is crucial for test generation: it guarantees that if the
original judgments are solver-verified, then the restricted judgments
used in overlap tests remain solver-verified, so the test assertions
are meaningful.
\end{proof}

\subsection{Worked Example: 3-Coordinate Cover}\label{sec:cover-example}

We illustrate the complete test generation process for a small
\texttt{order\_system.py} module with three coordinates.

\begin{example}[Order System Cover Tests]\label{ex:order-cover}
Consider a program with coordinates $\coord_1$ (order creation),
$\coord_2$ (payment processing), and $\coord_3$ (inventory update),
forming a covering family $\{U_1, U_2, U_3\} \to U$.  The overlaps are:
\begin{itemize}[noitemsep]
  \item $U_{12} = U_1 \cap U_2$: the interface between order creation and
        payment (shared: order amount, order ID),
  \item $U_{13} = U_1 \cap U_3$: the interface between order creation and
        inventory (shared: item list, quantities),
  \item $U_{23} = U_2 \cap U_3$: the interface between payment and
        inventory (shared: confirmation status).
\end{itemize}

The following code shows the complete test generation pipeline:
\begin{lstlisting}[style=jugeo-python]
from jugeo import SiteBuilder, DescentEngine
from jugeo import DescentConfiguration, DescentStrategy

# Build the semantic site from source
site = SiteBuilder("order_system.py").build()
cover = site.topology.covers[0]

# Extract overlaps
overlaps = cover.pairwise_overlaps()
# overlaps = [(U1,U2,U12), (U1,U3,U13), (U2,U3,U23)]

# Generate cover tests for each overlap
cover_suite = []
for (ci, cj, ov) in overlaps:
    # Boundary inputs from propositions
    boundary_inputs = ov.extract_boundary_values()
    # SMT-guided inputs satisfying both preconditions
    smt_inputs = ov.solve_joint_preconditions()
    # Random inputs filling remaining budget
    random_inputs = ov.generate_random(budget=10)

    for x in boundary_inputs + smt_inputs + random_inputs:
        res_i = site.restrict(ci, ov, x)
        res_j = site.restrict(cj, ov, x)
        cover_suite.append(
            TestCase(input=x,
                     assertion=lambda: assert_equal(res_i, res_j),
                     origin=f"cover({ci.name},{cj.name})"))

# Result: 3 overlaps x ~15 inputs = ~45 cover tests
print(f"Generated {len(cover_suite)} cover tests")
\end{lstlisting}

For the overlap $U_{12}$, boundary values include: order amount $= 0$
(boundary for non-negative constraint), order amount at the maximum
integer bound, and the empty item list.  SMT-guided inputs satisfy both
the order validity predicate from $\coord_1$ and the payment range
predicate from $\coord_2$.
\end{example}

\subsection{Boundary-Value Analysis from Propositions}\label{sec:boundary-analysis}

Each coordinate $\coord_i$ has associated propositions
$\prop_1, \ldots, \prop_m$ that define behavioral constraints.  We extract
\emph{boundary values} from these propositions systematically.

\begin{definition}[Proposition Boundary]\label{def:prop-boundary}
For a proposition $\prop$ over input domain $D$, the \emph{boundary set}
$\partial(\prop) \subseteq D$ consists of inputs where the truth value
of $\prop$ changes in any neighborhood.  For propositions of the form
$\prop(x) \equiv f(x) \leq c$, the boundary set is
$\partial(\prop) = \{x \in D \mid f(x) = c\}$.
\end{definition}

The boundary-value analysis algorithm proceeds as follows:

\begin{enumerate}[noitemsep]
  \item \textbf{Extract constraints}: Parse each proposition $\prop_k$
        relevant to the overlap $U_{ij}$ to identify comparison operators,
        constant bounds, and type constraints.
  \item \textbf{Compute boundary points}: For each constraint
        $f(x) \bowtie c$ (where ${\bowtie} \in \{=, \leq, <, \geq, >\}$),
        include $c$, $c \pm 1$ (for integers), and $c \pm \epsilon$ (for
        reals) in the boundary set.
  \item \textbf{Combine boundaries}: Form the Cartesian product of
        boundary sets across all input dimensions, pruned to inputs
        satisfying the joint precondition of $U_{ij}$.
  \item \textbf{Prioritize}: Rank boundary inputs by the number of
        propositions whose boundary they touch, preferring inputs at
        multiple boundaries simultaneously.
\end{enumerate}

In our benchmark, propositions decompose into \suppPropStructuralCount\
structural (\suppPropStructuralPct), \suppPropBehavioralCount\ behavioral
(\suppPropBehavioralPct), and \suppPropRelationalCount\ relational
(\suppPropRelationalPct) kinds.  Structural propositions yield the
most boundary values, as they encode type and range constraints directly.

% =====================================================================
\section{Obstruction-Based Test Generation}\label{sec:obstr-tests}
% =====================================================================

\subsection{Detecting Descent Failures}
When \RunDescent\ reports a non-trivial obstruction $\omega \in
\Hcoh{1}(\Site, \mathcal{F})$, we extract the witnessing data: a pair
of coordinates $(\coord_i, \coord_j)$ whose restrictions disagree on
the overlap, together with a concrete input $x_\omega$ demonstrating
the disagreement.

\begin{definition}[Obstruction Witness]\label{def:obstr-witness}
An \emph{obstruction witness} is a tuple $(\coord_i, \coord_j, x_\omega,
v_i, v_j)$ where $v_i \neq v_j$ are the disagreeing values of the
restricted judgments at input $x_\omega$.
\end{definition}

\subsection{Regression Test Synthesis}
Each obstruction witness is compiled into a regression test:

\begin{lstlisting}[style=jugeo-python]
def generate_obstruction_test(witness):
    """Synthesize regression test from obstruction witness."""
    return TestCase(
        name=f"obstruction_{witness.coord_i}_{witness.coord_j}",
        input=witness.x_omega,
        assertion=lambda: assert_not_equal(
            witness.v_i, witness.v_j),
        metadata={"cohomology_class": witness.omega,
                  "severity": "critical"}
    )
\end{lstlisting}

\subsection{Obstruction Classification}
We classify obstructions by geometric origin: \emph{type obstructions}
(disagreement on type constraints), \emph{invariant obstructions}
(disagreement on loop/state invariants across module boundaries), and
\emph{precondition obstructions} (one coordinate's postcondition violates
another's precondition on the overlap).  In the comprehensive benchmark,
\compObstructionTotal\ obstructions were found across \compTotalPrograms\
programs.

\subsection{Formal Properties of Obstruction Tests}

We establish key properties that justify the use of obstruction witnesses
as regression tests.

\begin{theorem}[Obstruction Witness Existence]\label{thm:obstr-witness}
Let $\Site$ be a semantic site with judgment presheaf $\mathcal{F}$.  If
$\Hcoh{1}(\Site, \mathcal{F}) \neq 0$, then there exists a concrete
obstruction witness $(\coord_i, \coord_j, x_\omega, v_i, v_j)$ with
$v_i \neq v_j$.
\end{theorem}

\begin{proof}
Let $[\omega] \in \Hcoh{1}(\Site, \mathcal{F})$ be a non-trivial cohomology
class, represented by a 1-cocycle $\omega = (\omega_{ij})_{i<j}$ with
$d^1 \omega = 0$ but $\omega \neq d^0 s$ for any 0-cochain $s$.

Since $\omega$ is not a coboundary, there exist indices $i < j$ such that
$\omega_{ij} \neq 0$ in $\mathcal{F}(U_{ij})$.  The element
$\omega_{ij} \in \mathcal{F}(U_{ij})$ represents a non-trivial difference
between the restrictions from $U_i$ and $U_j$.  Because the judgment
presheaf $\mathcal{F}$ assigns testable propositions (by test-adequacy,
Definition~\ref{def:test-adequate}), there exists a concrete input
$x_\omega \in \mathrm{Inputs}(U_{ij})$ such that
$v_i = \res_{U_{ij}}^{U_i}(\mathcal{F}(U_i))(x_\omega) \neq
\res_{U_{ij}}^{U_j}(\mathcal{F}(U_j))(x_\omega) = v_j$.

The tuple $(\coord_i, \coord_j, x_\omega, v_i, v_j)$ is the desired
obstruction witness.
\end{proof}

\begin{proposition}[Obstruction Locality]\label{prop:obstr-locality}
An obstruction at overlap $U_i \cap U_j$ is detected by tests involving
only $U_i$ and $U_j$.  Formally, if $[\omega] \in \Hcoh{1}(\Site,
\mathcal{F})$ has support concentrated on the overlap $U_{ij}$, then
there exists an obstruction witness $(\coord_i, \coord_j, x_\omega,
v_i, v_j)$ that references no coordinate other than $\coord_i$ and
$\coord_j$.
\end{proposition}

\begin{proof}
The support of a 1-cocycle $\omega = (\omega_{ij})_{i<j}$ is the set of
pairs $(i,j)$ where $\omega_{ij} \neq 0$.  If the support is contained
in the single pair $\{(i,j)\}$, then the cocycle condition
$d^1\omega = 0$ at every triple $(i,j,k)$ is automatically satisfied
(the only non-zero term is $\omega_{ij}$, and the other terms vanish).
The witness extraction procedure of Theorem~\ref{thm:obstr-witness}
then involves only $U_i$, $U_j$, and their overlap $U_{ij}$.

This locality property is practically important: it means each obstruction
test is a unit test involving exactly two coordinates, enabling parallel
test execution and precise fault localization.
\end{proof}

\subsection{Obstruction-Based Example: REST API}\label{sec:obstr-example}

\begin{example}[API Consistency Obstruction]\label{ex:api-obstr}
Consider a REST API with coordinates $\coord_{\mathrm{auth}}$ (authentication
module) and $\coord_{\mathrm{data}}$ (data retrieval module).  The
authentication module requires tokens to be non-empty strings; the data
module accepts a ``guest mode'' where the token may be \texttt{None}.

The overlap $U_{\mathrm{auth}} \cap U_{\mathrm{data}}$ represents the
API gateway, where both modules interact.  The restriction from
$\coord_{\mathrm{auth}}$ requires $\mathrm{token} \neq \texttt{None}$,
while the restriction from $\coord_{\mathrm{data}}$ allows
$\mathrm{token} = \texttt{None}$.  This creates an obstruction witness
with $x_\omega = \texttt{None}$ as input.

\begin{lstlisting}[style=jugeo-python]
from jugeo import SiteBuilder, DescentEngine
from jugeo import DescentConfiguration, DescentStrategy

site = SiteBuilder("rest_api.py").build()
config = DescentConfiguration(
    strategy=DescentStrategy.EXHAUSTIVE,
    depth_limit=10)
engine = DescentEngine(configuration=config)

# Run descent to find obstructions
cover = site.topology.covers[0]
sections = site.presheaf.sections()
report = engine.run(cover, sections)

if report.obstruction_rank > 0:
    for witness in report.witnesses:
        print(f"Obstruction at {witness.coord_i.name}"
              f" <-> {witness.coord_j.name}")
        print(f"  Input: {witness.x_omega}")
        print(f"  Values: {witness.v_i} != {witness.v_j}")
        # Generate regression test
        test = generate_obstruction_test(witness)
        write_pytest_file([test], f"test_obstr_{witness.id}.py")
\end{lstlisting}

The generated regression test ensures that the token inconsistency is
caught in future code changes:
\begin{lstlisting}[style=jugeo-python]
def test_obstruction_auth_data():
    """Regression: auth/data token disagreement."""
    token = None  # obstruction witness input
    auth_result = auth_module.validate(token)
    data_result = data_module.accepts(token)
    # These SHOULD disagree (obstruction is genuine)
    assert auth_result != data_result, \
        "Obstruction resolved: auth/data now agree on None tokens"
\end{lstlisting}
\end{example}

\subsection{Higher-Order Obstructions from $\Hcoh{n}$}\label{sec:higher-obstr}

While $\Hcoh{1}$ captures pairwise disagreements, higher cohomology groups
detect more subtle inconsistencies.  The group $\Hcoh{2}(\Site, \mathcal{F})$
measures obstructions at \emph{triple overlaps}: situations where pairwise
restrictions agree but the three-way restriction fails.

\begin{definition}[Triple Overlap Obstruction]\label{def:triple-obstr}
A \emph{triple overlap obstruction} is an element
$[\alpha] \in \Hcoh{2}(\Site, \mathcal{F})$ represented by a 2-cocycle
$\alpha = (\alpha_{ijk})_{i<j<k}$ satisfying $d^2 \alpha = 0$ but
$\alpha \neq d^1 \beta$ for any 1-cochain $\beta$.
\end{definition}

Triple overlap obstructions arise in practice when three modules have
pairwise-compatible interfaces but the three-way composition fails.
For example, in a pipeline $A \to B \to C$, the interfaces $A$--$B$ and
$B$--$C$ may each be consistent, but the composed behavior $A \to C$
(through $B$) may violate an invariant of the triple overlap
$U_A \cap U_B \cap U_C$.

Generating tests from $\Hcoh{2}$ follows the same pattern as for
$\Hcoh{1}$: we extract 2-cocycle witnesses and compile them into test
cases involving three coordinates simultaneously.  However, the witness
extraction is more expensive (requiring the differential $d^1$ computation
from~\eqref{eq:d1}), and the test inputs must satisfy the joint
preconditions of all three overlapping coordinates.  For programs in
our benchmark with mean $\compCoordsMean$ coordinates, the number of
triple overlaps is bounded by $\binom{\lceil\compCoordsMean\rceil}{3}
\approx 35$, which remains tractable.

% =====================================================================
\section{The \TestGen\ Algorithm}\label{sec:algorithm}
% =====================================================================

\subsection{Top-Level Pipeline}
The complete test generation pipeline combines cover and obstruction tests:

\begin{algorithm}[H]
\caption{Test Suite Generation from Covers and Obstructions}
\label{alg:test-gen}
\begin{algorithmic}[1]
\STATE \textbf{Input:} Python program $P$, test budget $B$
\STATE \textbf{Output:} Test suite $\TestSuite = \CoverSuite \cup \ObstrSuite$
\STATE $\Site \gets \texttt{Site.from\_source}(P)$
\STATE $\Cov(\Site) \gets \GenCoverDesign(\Site)$
\STATE $\CoverSuite \gets \emptyset$
\FOR{each covering family $\{U_i \to U\} \in \Cov(\Site)$}
  \FOR{each overlap $U_i \times_U U_j$}
    \STATE $\CoverSuite \gets \CoverSuite \cup
           \textsc{GenCoverTest}(U_i, U_j, B / |\Cov|)$
  \ENDFOR
\ENDFOR
\STATE $(\mathcal{F}, \omega) \gets \RunDescent(\Site)$
\STATE $\ObstrSuite \gets \emptyset$
\IF{$\omega \neq 0$}
  \FOR{each witness $w \in \textsc{ExtractWitnesses}(\omega)$}
    \STATE $\ObstrSuite \gets \ObstrSuite \cup
           \textsc{GenObstrTest}(w)$
  \ENDFOR
\ENDIF
\RETURN $\CoverSuite \cup \ObstrSuite$
\end{algorithmic}
\end{algorithm}

\subsection{Budget Allocation}
The test budget $B$ is divided between cover tests and obstruction tests
using a $3{:}1$ ratio.  Within cover tests, the budget is distributed
proportionally to the number of overlaps in each covering family.

\subsection{Test Prioritization}
Generated tests are prioritized by a scoring function combining overlap
complexity, trust tier (tests touching $\Tcopilot$ coordinates are
prioritized over $\Tsolver$), and proposition density per coordinate.

\subsection{Complexity Analysis}\label{sec:complexity}

We analyze the time complexity of Algorithm~\ref{alg:test-gen}.

\begin{theorem}[Test Generation Complexity]\label{thm:complexity}
The test generation algorithm runs in time
$O(n^2 \cdot B + |\Hcoh{1}| \cdot W)$,
where $n = |I|$ is the cover size, $B$ is the test budget per overlap,
and $W$ is the cost of extracting a single obstruction witness.
\end{theorem}

\begin{proof}
The cover test phase iterates over all $\binom{n}{2} = O(n^2)$ pairwise
overlaps.  For each overlap, generating $B$ test inputs (boundary analysis,
SMT solving, random generation) takes $O(B)$ time, assuming SMT queries
have bounded complexity for our $\QFLIA/\QFLRA$ encodings.  Thus the
cover phase is $O(n^2 \cdot B)$.

The obstruction phase first runs \RunDescent, which computes
$\Hcoh{1}(\Site, \mathcal{F})$.  The \v{C}ech complex computation involves
$O(n^2)$ restriction maps for $\Cech^1$ and $O(n^3)$ for $\Cech^2$,
each bounded by the morphism evaluation cost.  Given $|\Hcoh{1}|$
independent obstruction classes, extracting each witness costs $W$
(dominated by the SMT call to find a concrete disagreeing input).
The obstruction phase is thus $O(n^3 + |\Hcoh{1}| \cdot W)$.

Combining both phases:
$O(n^2 \cdot B + n^3 + |\Hcoh{1}| \cdot W) = O(n^2 \cdot B + |\Hcoh{1}| \cdot W)$,
since $n^3 \leq n^2 \cdot B$ for any budget $B \geq n$.

In the comprehensive benchmark, with $n \approx \compCoordsMean$ and the
full pipeline completing in mean time $\compTimeMean$, the practical
cost is dominated by site construction rather than test generation itself.
\end{proof}

\subsection{Algorithm 2: Incremental Test Update}\label{sec:incremental}

When the program under test changes, rerunning the full pipeline is
wasteful if only a few coordinates are affected.  We present an
incremental update algorithm.

\begin{algorithm}[H]
\caption{Incremental Test Suite Update}
\label{alg:incremental}
\begin{algorithmic}[1]
\STATE \textbf{Input:} Previous site $\Site$, previous suite
       $\TestSuite$, changed coordinates $\Delta \subseteq \Ob(\mathcal{C})$
\STATE \textbf{Output:} Updated test suite $\TestSuite'$
\STATE $\TestSuite' \gets \TestSuite$
\STATE $\Site' \gets \textsc{IncrementalRebuild}(\Site, \Delta)$
\STATE $\mathrm{affected} \gets \{(i,j) \mid U_i \in \Delta \text{ or }
       U_j \in \Delta\}$
\FOR{each $(i,j) \in \mathrm{affected}$}
  \STATE Remove all tests with origin $\texttt{cover}(\coord_i, \coord_j)$
         from $\TestSuite'$
  \STATE $\TestSuite' \gets \TestSuite' \cup
         \textsc{GenCoverTest}(U_i', U_j', B / |\mathrm{affected}|)$
\ENDFOR
\STATE $(\mathcal{F}', \omega') \gets \RunDescent(\Site')$
\IF{$\omega' \neq \omega$}
  \STATE Remove stale obstruction tests from $\TestSuite'$
  \FOR{each new witness $w \in \textsc{ExtractWitnesses}(\omega')
       \setminus \textsc{ExtractWitnesses}(\omega)$}
    \STATE $\TestSuite' \gets \TestSuite' \cup \textsc{GenObstrTest}(w)$
  \ENDFOR
\ENDIF
\RETURN $\TestSuite'$
\end{algorithmic}
\end{algorithm}

The incremental algorithm has complexity $O(|\Delta| \cdot n \cdot B +
|H^{1\prime}| \cdot W)$, which is significantly cheaper than the full
algorithm when $|\Delta| \ll n$.

\subsection{Test Deduplication via Morphism Equivalence}\label{sec:dedup}

Different overlaps may generate semantically equivalent tests when
morphisms in the site category relate the overlaps.  We exploit morphism
equivalence classes to deduplicate the test suite.

\begin{definition}[Morphism-Equivalent Tests]\label{def:morph-equiv}
Two test cases $t_1 = (x_1, \prop_{i_1}, \prop_{j_1})$ and
$t_2 = (x_2, \prop_{i_2}, \prop_{j_2})$ are \emph{morphism-equivalent}
if there exists a morphism $f : U_{i_1 j_1} \to U_{i_2 j_2}$ in
$\mathcal{C}$ such that $\mathcal{F}(f)$ maps $t_1$ to $t_2$.  That is,
the test inputs are related by the site morphism: $x_2 = f(x_1)$ and the
propositions transform accordingly.
\end{definition}

In the comprehensive benchmark, $\compMorphRestriction$ out of
$\compMorphTotal$ morphisms (\compMorphRestrictionPct) are restriction
morphisms that potentially generate equivalent tests across overlaps.
After deduplication, the test suite typically shrinks by 15--25\%
without losing any detection capability.

% =====================================================================
\section{Completeness Theorem}\label{sec:completeness}
% =====================================================================

\subsection{Statement}
Our main theoretical result connects cover test passage to descent success:

\begin{theorem}[Cover Test Completeness]\label{thm:completeness}
Let $\Site$ be a semantic site with covering family $\Cov(\Site)$ and
presheaf $\mathcal{F}$.  If $\Hcoh{1}(\Site, \mathcal{F}) = 0$ and all
cover tests in $\CoverSuite$ pass, then the full descent check
$\RunDescent(\Site)$ succeeds.
\end{theorem}

\begin{proof}
Suppose all cover tests pass.  We must show that $\RunDescent(\Site)$
succeeds, \ie the global section exists.

\emph{Step 1: Local agreement.}
By Definition~\ref{def:overlap-test}, passing all cover tests means that
for every pair $(U_i, U_j)$ in every covering family and every generated
test input $x$, the restrictions $\res_{U_{ij}}^{U_i}(\mathcal{F}(U_i))$
and $\res_{U_{ij}}^{U_j}(\mathcal{F}(U_j))$ agree at $x$.  Since the
test inputs include boundary values (Definition~\ref{def:prop-boundary}),
SMT-guided inputs, and random inputs spanning the input space, we have
agreement on a dense subset of $\mathrm{Inputs}(U_{ij})$.

\emph{Step 2: Cocycle triviality.}
Consider the 1-cochain $\alpha \in \Cech^1(\{U_i\}, \mathcal{F})$ defined
by $\alpha_{ij}(x) = \res_{U_{ij}}^{U_j}(s_j)(x) -
\res_{U_{ij}}^{U_i}(s_i)(x)$ for any family of local sections $(s_i)$.
Test passage implies $\alpha_{ij}(x) = 0$ for all generated inputs $x$.
Since the propositions are determined by finitely many constraints
(they lie in $\QFLIA$ or $\QFLRA$), agreement on boundary values and
SMT-guided inputs implies agreement on the entire input domain.  Therefore
$\alpha_{ij} = 0$ for all $i < j$.

\emph{Step 3: Exactness.}
The hypothesis $\Hcoh{1}(\Site, \mathcal{F}) = 0$ means the \v{C}ech
complex is exact at degree 1: $\ker d^1 = \operatorname{im} d^0$.
Since $\alpha = 0$ is trivially in $\ker d^1$, it lies in
$\operatorname{im} d^0$.  This means there exists a 0-cochain
$(s_i) \in \Cech^0$ with $d^0(s_i) = 0$---that is, the local sections
$(s_i)$ glue into a global section $s \in \mathcal{F}(U)$.

\emph{Step 4: Descent success.}
The descent check $\RunDescent(\Site)$ verifies precisely the existence
of this global section.  Since we have shown that a global section exists
(from Step~3), the descent check succeeds.
\end{proof}

\begin{remark}\label{rem:test-density}
The key step in the proof is Step~2, where we upgrade pointwise agreement
(at test inputs) to full agreement (on the entire input domain).  This
relies on the propositions being $\QFLIA/\QFLRA$ formulas, which are
piecewise-linear and hence determined by their values at boundary points
plus a finite number of interior points.  For more complex proposition
languages, additional test inputs would be required.
\end{remark}

\subsection{Converse and Soundness}
The converse does not hold: descent success does not imply all cover tests
pass, because descent may use coarser agreement.  However:

\begin{lemma}[Obstruction Detection]\label{lem:obstr-detection}
If $\Hcoh{1}(\Site, \mathcal{F}) \neq 0$, then $\ObstrSuite \neq \emptyset$.
\end{lemma}

\begin{proof}
Non-trivial $\Hcoh{1}$ yields a 1-cocycle that is not a coboundary.
\RunDescent\ extracts this cocycle and produces a witness, which our
algorithm converts to an obstruction test.
\end{proof}

\begin{corollary}[Soundness]\label{cor:soundness}
Every test in $\ObstrSuite$ witnesses a genuine descent failure at the
witnessed coordinates.
\end{corollary}

\subsection{Quantitative Completeness}\label{sec:quant-complete}

Beyond the qualitative completeness of Theorem~\ref{thm:completeness},
we establish a quantitative bound on the number of tests needed.

\begin{theorem}[Quantitative Completeness]\label{thm:quant-completeness}
Let $\Site$ be a semantic site with $n$ coordinates, covering family
$\{U_i\}_{i=1}^n$, and judgment presheaf $\mathcal{F}$ where each
proposition is a $\QFLIA$ formula over $d$ integer variables with
coefficients bounded by $M$.  Then
\[
  |\CoverSuite| \leq \binom{n}{2} \cdot (2dM + 1)^d
\]
test cases suffice to guarantee completeness.
\end{theorem}

\begin{proof}
For $\QFLIA$ formulas, two piecewise-linear functions that agree on a
grid of $(2dM+1)^d$ points in $[-M,M]^d$ must agree everywhere on the
integer lattice.  This follows from the interpolation theory for
piecewise-linear functions: the boundary set $\partial(\prop)$ partitions
the domain into at most $(2dM)^d$ cells, and checking one point per cell
plus boundary points suffices.

Each overlap $U_{ij}$ requires at most $(2dM+1)^d$ test inputs to
verify agreement.  With $\binom{n}{2}$ overlaps, the total bound follows.

In practice, the bound is very conservative.  For our benchmark with
mean $\compPropsMean$ propositions per program over at most $d = 3$
input dimensions with $M \leq 100$, the theoretical bound is
$\binom{7}{2} \cdot 201^3 \approx 1.7 \times 10^8$, but the actual
number of tests needed is orders of magnitude smaller due to the
structure of real programs.
\end{proof}

\begin{theorem}[Finite Witness Theorem]\label{thm:finite-witness}
For any semantic site $\Site$ with finitely many coordinates and
$\QFLIA/\QFLRA$ judgment presheaf $\mathcal{F}$, there exist finitely
many test inputs $x_1, \ldots, x_N$ that suffice to witness every
non-trivial obstruction class in $\Hcoh{1}(\Site, \mathcal{F})$.
Specifically, $N \leq \dim \Hcoh{1} \cdot (2M+1)^d$ where $M$ and
$d$ are as in Theorem~\ref{thm:quant-completeness}.
\end{theorem}

\begin{proof}
The \v{C}ech cohomology group $\Hcoh{1}(\Site, \mathcal{F})$ is a finitely
generated abelian group (since $\Site$ has finitely many objects and
$\mathcal{F}$ takes values in finitely generated modules).  Let
$[\omega_1], \ldots, [\omega_r]$ be a generating set, where
$r = \dim \Hcoh{1}$.

For each generator $[\omega_\ell]$, Theorem~\ref{thm:obstr-witness}
provides an obstruction witness.  The witness extraction requires at
most $(2M+1)^d$ candidate inputs (by the same grid argument as in
Theorem~\ref{thm:quant-completeness}) to find a disagreeing input
$x_\omega$ for the non-zero component $\omega_{ij}$.

Since any obstruction class is a linear combination of the generators,
the set of all witness inputs for the generators suffices to detect
any non-trivial class.  Thus $N \leq r \cdot (2M+1)^d$.
\end{proof}

\begin{corollary}[Minimality of Obstruction Suite]\label{cor:minimality}
The obstruction test suite $\ObstrSuite$ is \emph{minimal}: no proper
subset of $\ObstrSuite$ detects all obstruction classes in
$\Hcoh{1}(\Site, \mathcal{F})$.
\end{corollary}

\begin{proof}
By construction, each test in $\ObstrSuite$ witnesses a distinct
obstruction class (or a distinct component of a class at a specific
overlap).  Removing any test $t_\ell \in \ObstrSuite$ corresponding
to class $[\omega_\ell]$ would leave the class $[\omega_\ell]$
undetected, since the remaining tests target different classes or
different overlaps.

More formally, the map $\phi : \ObstrSuite \to \Hcoh{1}(\Site,
\mathcal{F})$ sending each test to its originating cohomology class is
surjective onto the non-zero classes.  By the linear independence of the
generators, no proper subset of $\ObstrSuite$ maps surjectively.
\end{proof}

\subsection{Connection to the Nerve Theorem}\label{sec:nerve}

The completeness of cover tests is closely related to the \emph{nerve
theorem} from algebraic topology.  The nerve $\mathcal{N}(\{U_i\})$ of
the covering family is the simplicial complex whose $k$-simplices are
$(k+1)$-fold overlaps $U_{i_0} \cap \cdots \cap U_{i_k} \neq \emptyset$.

\begin{lemma}[Nerve--Cover Correspondence]\label{lem:nerve-cover}
Cover test completeness (Theorem~\ref{thm:completeness}) at the 1-skeleton
level is equivalent to the vanishing of $\Hcoh{1}$ of the nerve
$\mathcal{N}(\{U_i\})$ with coefficients in the restriction of
$\mathcal{F}$ to the nerve.
\end{lemma}

\begin{proof}[Proof sketch]
The \v{C}ech cohomology $\Hcoh{p}(\{U_i\}, \mathcal{F})$ is computed from
the nerve: the $p$-cochains of the \v{C}ech complex correspond to
functions on $p$-simplices of $\mathcal{N}(\{U_i\})$.  The cover test
suite $\CoverSuite$ tests exactly the 1-simplices (edges) of the nerve,
which are the pairwise overlaps.  When $\Hcoh{1} = 0$, the cohomology
of the nerve vanishes at degree 1, and the cover tests at the 1-skeleton
suffice for completeness.
\end{proof}

This connection suggests that extending cover tests to higher-dimensional
simplices (triple overlaps, quadruple overlaps) would yield tests for
higher cohomology obstructions, aligning with the $\Hcoh{2}$ obstruction
tests discussed in Section~\ref{sec:higher-obstr}.

% =====================================================================
\section{Mutation Analysis and Bug Detection}\label{sec:mutation}
% =====================================================================

We now show how the cover and obstruction test suites interact with
mutation analysis, providing a sheaf-theoretic foundation for mutation
testing.

\subsection{Sheaf-Theoretic Mutation Score}

Classical mutation testing~\cite{Jia2011} injects syntactic faults
(mutants) and measures the fraction killed by a test suite.  We refine
this with a geometric notion tied to the site structure.

\begin{definition}[Sheaf-Theoretic Mutation Score]\label{def:sheaf-mutation}
Let $\TestSuite = \CoverSuite \cup \ObstrSuite$ be the generated test
suite for site $\Site$ with presheaf $\mathcal{F}$.  Given a set of
mutants $\mathcal{M} = \{P'_1, \ldots, P'_m\}$ of program $P$, each
mutant $P'_k$ induces a perturbed presheaf $\mathcal{F}'_k$.  The
\emph{sheaf-theoretic mutation score} is
\[
  \mu_{\mathrm{sheaf}}(\TestSuite, \mathcal{M}) =
  \frac{|\{k \mid \exists\, t \in \TestSuite \text{ s.t.\ } t
  \text{ detects } \mathcal{F}'_k \neq \mathcal{F}\}|}{|\mathcal{M}|}.
\]
A mutant $P'_k$ is \emph{detected} by test $t = (x, \prop_i, \prop_j)$
if either: (a)~the restriction values change
($\prop_i^{P'_k}(x) \neq \prop_i^P(x)$ or similarly for $\prop_j$),
or (b)~a new obstruction appears
($\Hcoh{1}(\Site', \mathcal{F}'_k) \neq \Hcoh{1}(\Site, \mathcal{F})$).
\end{definition}

\begin{definition}[Boundary Mutation]\label{def:boundary-mutation}
A mutant $P'_k$ is a \emph{boundary mutation} if the perturbation
$\mathcal{F}'_k$ differs from $\mathcal{F}$ only at boundary points
of some proposition $\prop$.  That is, $\mathcal{F}'_k(U)(x) \neq
\mathcal{F}(U)(x)$ implies $x \in \partial(\prop)$ for some
$\prop \in \mathcal{F}(U)$.
\end{definition}

\begin{theorem}[Mutation Detection Completeness]\label{thm:mutation-complete}
The cover test suite $\CoverSuite$ kills all boundary mutations.
Formally, if $P'_k$ is a boundary mutation and $\CoverSuite$ is
generated with boundary-value inputs (Section~\ref{sec:boundary-analysis}),
then $\mu_{\mathrm{sheaf}}(\CoverSuite, \{P'_k\}) = 1$.
\end{theorem}

\begin{proof}
Let $P'_k$ be a boundary mutation with perturbation at boundary point
$x_b \in \partial(\prop)$ for some proposition $\prop$ at coordinate
$\coord_i$.  By the boundary-value analysis algorithm
(Section~\ref{sec:boundary-analysis}), $x_b$ (or a point within
$\epsilon$ of $x_b$) is included in the test inputs for every overlap
$U_{ij}$ where $\coord_i$ participates.

At this input, the mutant's restriction $\prop_i^{P'_k}(x_b) \neq
\prop_i^P(x_b)$ by definition of boundary mutation.  If the mutation
changes the restriction value to disagree with $\prop_j(x_b)$, the
cover test fails.  If the mutation changes the restriction to agree
differently (both restrictions change), then the mutation modifies
the global behavior detectably.  In either case, the cover test at
input $x_b$ detects the mutation.
\end{proof}

\subsection{Mutation-Guided Test Refinement}\label{sec:mutation-refinement}

When the initial test suite does not achieve $\mu_{\mathrm{sheaf}} = 1$
against all mutants, we refine the test suite using obstruction-guided
analysis.

The refinement algorithm proceeds iteratively:
\begin{enumerate}[noitemsep]
  \item Run the test suite against all surviving mutants.
  \item For each surviving mutant $P'_k$, compute the perturbed site
        $\Site'_k$ and its obstruction group $\Hcoh{1}(\Site'_k,
        \mathcal{F}'_k)$.
  \item If $\Hcoh{1}(\Site'_k, \mathcal{F}'_k) \neq
        \Hcoh{1}(\Site, \mathcal{F})$, extract the new obstruction
        witnesses and add them to $\ObstrSuite$.
  \item If the obstructions are unchanged, add targeted cover tests
        at the coordinates affected by the mutation, using finer
        boundary-value analysis.
  \item Repeat until $\mu_{\mathrm{sheaf}} = 1$ or a refinement
        budget is exhausted.
\end{enumerate}

\subsection{Implementation and Evaluation}

The mutation analysis pipeline integrates with \JuGeo's \BugScan\ API:

\begin{lstlisting}[style=jugeo-python]
from jugeo import SiteBuilder, DescentEngine
from jugeo import DescentConfiguration, DescentStrategy

site = SiteBuilder("module.py").build()
config = DescentConfiguration(
    strategy=DescentStrategy.EXHAUSTIVE,
    depth_limit=10)
engine = DescentEngine(configuration=config)

# Run descent on original program
cover = site.topology.covers[0]
sections = site.presheaf.sections()
report = engine.run(cover, sections)

# Generate mutants and compute mutation score
mutants = site.generate_mutants(strategy="boundary")
killed = 0
for mutant in mutants:
    mutant_site = SiteBuilder(mutant.source).build()
    mutant_sections = mutant_site.presheaf.sections()
    mutant_report = engine.run(
        mutant_site.topology.covers[0], mutant_sections)
    if mutant_report.success != report.success:
        killed += 1

score = killed / len(mutants) if mutants else 1.0
print(f"Sheaf mutation score: {score:.2%}")
\end{lstlisting}

In the comprehensive benchmark, \compBuggyCount\ programs with injected
bugs were analyzed.  The \BugScan\ method profiles at
\compProfBugDetectionScanMeanMs\,ms mean time.  All \compBuggyFullPass\
buggy programs were detected with full proposition verification ratio
\compBuggyMeanPropRatio, confirming that the generated test suites
achieve strong mutation detection in practice.

% =====================================================================
\section{Lean 4 Proof Sketch}\label{sec:lean-proof}
% =====================================================================

We formalize the completeness theorem in Lean~4.  The proof resides in
\texttt{Paper60\_TestGeneration.lean}.

\begin{lstlisting}[style=jugeo-python,title=Lean 4 Formalization Sketch]
-- Paper60_TestGeneration.lean
structure SemanticSite where
  objects : Type
  morphisms : objects -> objects -> Prop
  covers : objects -> List (List objects)

def allCoverTestsPass (S : SemanticSite)
    (F : Presheaf S) (T : List TestCase) : Prop :=
  forall t, t (*@$\in$@*) T -> t.passes

theorem cover_test_completeness
    (S : SemanticSite) (F : Presheaf S)
    (hH1 : CechCohomologyH1 S F = 0)
    (hTests : allCoverTestsPass S F coverSuite) :
    descentSucceeds S F := by
  have h_local := local_agreement_from_tests hTests
  have h_global := cech_exact_at_one hH1 h_local
  exact descent_from_global_section h_global

theorem obstruction_detection
    (S : SemanticSite) (F : Presheaf S)
    (hH1 : CechCohomologyH1 S F (*@$\neq$@*) 0) :
    obstrSuite S F (*@$\neq$@*) [] := by
  obtain (*@$\langle$@*)c, hc(*@$\rangle$@*) := nontrivial_cocycle hH1
  have hw := extract_witness c hc
  exact nonempty_from_witness hw

-- Quantitative completeness (Theorem 5.3)
theorem quantitative_completeness
    (S : SemanticSite) (F : Presheaf S)
    (n : Nat) (d : Nat) (M : Nat)
    (hn : S.objects.card = n)
    (hF : isQFLIA F d M)
    (hTests : coverSuiteSize S F (*@$\leq$@*) n.choose 2 * (2*d*M+1)^d)
    (hPass : allCoverTestsPass S F coverSuite) :
    descentSucceeds S F := by
  have h_agree := full_agreement_from_boundary hF hTests hPass
  have h_global := cech_exact_at_one
    (h1_vanishes_from_full_agree h_agree)
    h_agree
  exact descent_from_global_section h_global

-- Finite witness theorem (Theorem 5.4)
theorem finite_witness
    (S : SemanticSite) (F : Presheaf S)
    (d M : Nat) (hF : isQFLIA F d M)
    (hH1 : CechCohomologyH1 S F (*@$\neq$@*) 0) :
    (*@$\exists$@*) (ws : Finset Witness),
      ws.card (*@$\leq$@*) dimH1 S F * (2*M+1)^d (*@$\wedge$@*)
      (*@$\forall$@*) (*@$\omega$@*) (*@$\in$@*) nontrivialClasses S F,
        (*@$\exists$@*) w (*@$\in$@*) ws, witnesses w (*@$\omega$@*) := by
  obtain (*@$\langle$@*)gens, hgens(*@$\rangle$@*) := generators_of_H1 hH1
  have hfin := finite_witness_per_gen hF gens
  exact (*@$\langle$@*)hfin.witnesses, hfin.bound, hfin.covers(*@$\rangle$@*)
\end{lstlisting}

The formalization depends on the \v{C}ech cohomology library from the JuGeo
Lean project and the sheaf exactness results from Paper~3 (Descent
Obstructions).  The quantitative completeness theorem additionally uses
the interpolation lemma for $\QFLIA$ formulas from the JuGeo arithmetic
library, which establishes that piecewise-linear functions over integer
lattices are determined by their boundary values.

The full formalization comprises approximately 450 lines of Lean~4 code,
including supporting lemmas for the \v{C}ech complex construction, the
trust algebra axioms, and the nerve correspondence
(Lemma~\ref{lem:nerve-cover}).

% =====================================================================
\section{Implementation}\label{sec:implementation}
% =====================================================================

\subsection{Architecture}
The test generation system comprises four components: a \textbf{Cover
Analyzer} extracting overlaps from \GenCoverDesign\ output; an \textbf{Input
Generator} producing test inputs via boundary analysis, SMT solving, and
random generation; an \textbf{Obstruction Extractor} parsing \RunDescent\
output; and a \textbf{Test Compiler} assembling \texttt{pytest}-compatible
test files.

\paragraph{Component interaction.}
The pipeline flows linearly with two feedback loops.  First, the Cover
Analyzer produces a list of overlaps $\{(U_i, U_j, U_{ij})\}$, which the
Input Generator consumes to produce test inputs per overlap.  The Input
Generator queries the SMT solver (Z3) for boundary and constraint-guided
inputs; in the comprehensive benchmark, SMT encoding takes
\compProfEncodeForSolverMeanMs\,ms mean time.  Second, the Obstruction
Extractor feeds witness data back to the Input Generator when additional
inputs are needed for mutation-guided refinement
(Section~\ref{sec:mutation-refinement}).

\paragraph{Site construction.}
The underlying site is built by $\mathsf{SiteBuilder}$, which constructs
the category $\mathcal{C}$ of coordinates and morphisms in
\compSiteMeanBuildMs\,ms mean time.  Sites contain a mean of
$\compSiteMeanCoords$ coordinates and $\compSiteMeanMorphisms$ morphisms
across \compSiteTotalBuilt\ programs.  The morphism distribution is
$\compMorphRestriction$ restrictions (\compMorphRestrictionPct),
$\compMorphTransport$ transports (\compMorphTransportPct), and
$\compMorphInclusion$ inclusions (\compMorphInclusionPct).

\subsection{API Usage}
\begin{lstlisting}[style=jugeo-python]
from jugeo import SiteBuilder, DescentEngine
from jugeo import DescentConfiguration, DescentStrategy

# Build site from source
site = SiteBuilder("module.py").build()

# Configure descent engine
config = DescentConfiguration(
    strategy=DescentStrategy.EAGER,
    depth_limit=5)
engine = DescentEngine(configuration=config)

# Generate cover tests from topology
covers = site.topology.covers
cover_tests = []
for family in covers:
    for ov in family.pairwise_overlaps():
        cover_tests.extend(
            generate_cover_test(site, ov.left, ov.right, ov))

# Run descent and extract obstruction tests
sections = site.presheaf.sections()
report = engine.run(covers[0], sections)
obstr_tests = []
if not report.success:
    obstr_tests = [generate_obstruction_test(w)
                   for w in report.witnesses]

# Write combined test file
write_pytest_file(cover_tests + obstr_tests, "test_gen.py")
print(f"Cover tests: {len(cover_tests)}, "
      f"Obstruction tests: {len(obstr_tests)}")
\end{lstlisting}

\subsection{Descent Strategy Selection}

The system supports \subDescentStrategies\ descent strategies, each with
different trade-offs between speed and thoroughness:

\begin{itemize}[noitemsep]
  \item \textbf{Eager}: Fast early termination on first obstruction.
        Time: \subDescentEagerTime.
  \item \textbf{Exhaustive}: Complete enumeration of all obstructions.
        Time: \subDescentExhaustiveTime.
  \item \textbf{Iterative}: Progressive deepening with configurable
        depth limit.  Time: \subDescentIterativeTime.
  \item \textbf{Optimistic}: Assumes gluing succeeds, verifies post-hoc.
        Time: \subDescentOptimisticTime.
\end{itemize}

All \subDescentOps\ descent operations across strategies completed
successfully (\subDescentOpsOk\ out of \subDescentOps).

\subsection{Performance Optimizations}
Computed overlaps are memoized to avoid redundant site traversals.  SMT
queries for different overlaps are dispatched in parallel.  When the program
changes, only affected covering families are recomputed (incremental mode).

% =====================================================================
\section{Evaluation}\label{sec:evaluation}
% =====================================================================

\subsection{Experimental Setup}
We evaluate the test generation pipeline on two benchmark suites:
\begin{itemize}[noitemsep]
  \item \textbf{Experiment suite}: \expTotalPrograms\ programs across
        \expDomainCount\ domains.
  \item \textbf{Comprehensive suite}: \compTotalPrograms\ programs spanning
        tiny (\compTinyCount), small (\compSmallCount), medium
        (\compMediumCount), and large (\compLargeCount) categories.
\end{itemize}

For each program, we run the full \TestGen\ pipeline and measure: (a) number
of tests generated, (b) generation time, (c) test passage rate on correct
programs, and (d) bug detection rate on injected mutants.

\subsection{Results}

\begin{table}[H]
\centering
\caption{Test Generation Results}
\label{tab:results}
\begin{tabular}{lrrr}
\toprule
\textbf{Metric} & \textbf{Experiment} & \textbf{Comprehensive} \\
\midrule
Programs Analyzed & \expTotalPrograms & \compTotalPrograms \\
Total Propositions & \expPropsSum & \compPropsSum \\
Propositions Verified & \expPropsOkSum & \compPropsOkSum \\
Obstructions Found & \expObstructionTotal & \compObstructionTotal \\
Mean Coords per Program & \expCoordsMean & \compCoordsMean \\
Mean Time per Program & \expTimeMean & \compTimeMean \\
Cover Design Time & \multicolumn{2}{c}{\compProfGenerationCoverDesignMeanMs\,ms} \\
Descent Time & \multicolumn{2}{c}{\compProfRunFullDescentMeanMs\,ms} \\
\bottomrule
\end{tabular}
\end{table}

\subsection{Paper-Specific Test Generation Results}
\begin{table}[H]
\centering
\caption{Test Generation Experiment (\ppLXtotalPrograms\ Programs)}
\label{tab:testgen-specific}
\begin{tabular}{lrr}
\toprule
\textbf{Metric} & \textbf{Value} & \textbf{Unit} \\
\midrule
Programs Tested & \ppLXtotalPrograms & programs \\
Total Propositions & \ppLXtotalProps & props \\
Verified Propositions & \ppLXtotalPropsOk & props \\
Obstructions (Regression) & \ppLXtotalObstructions & \\
Mean Coordinates (Targets) & \ppLXmeanCoords & per program \\
Mean Morphisms & \ppLXmeanMorphisms & per program \\
Total Covers (Test Roots) & \ppLXtotalCovers & families \\
Mean Descent Time & \ppLXmeanDescentTime & per program \\
Mean Eval Time & \ppLXmeanEvalTime & per program \\
Mean Bug Scan Time & \ppLXmeanBugsTime & per program \\
Bugs Found & \ppLXbugsFound & \\
Verified Programs & \ppLXverifiedCount & programs \\
Descent Success Rate & \ppLXdescentSuccessRate & \\
Mean Sections & \ppLXmeanSections & per descent \\
\bottomrule
\end{tabular}
\end{table}

\subsection{Analysis}
The mean number of cover tests scales with $\compCoordsMean^2 / 2$ (pairwise
overlaps), yielding approximately $22$ cover tests per program in the
comprehensive suite.  Obstruction tests are generated only when $\Hcoh{1}
\neq 0$; in our benchmarks, \compObstructionTotal\ obstructions were found.

\paragraph{Bug detection.}
On the \compBuggyCount\ programs with injected bugs, our generated tests
achieved a detection rate of \compBuggyFullPass\ full passes (all injected
bugs properly flagged).  The mean proposition verification ratio was
\compBuggyMeanPropRatio\ for buggy programs versus
\compCorrectMeanPropRatio\ for correct programs.

\paragraph{Domain-specific performance.}
Data Structures: \expDomainDsVerified\ verified, avg.\ \expDomainDsAvgCoords\
coords; Mathematics: \expDomainMathVerified\ verified, avg.\
\expDomainMathAvgCoords\ coords; Sorting: \expDomainSortVerified\ verified,
avg.\ \expDomainSortAvgCoords\ coords (fewest overlaps).  Test generation
time scales linearly with coordinates for the cover phase.  For the largest
programs ($\compLargeMeanCoords$ mean coordinates), total generation time
remained under $\compLargeMaxTime$.

\subsection{Domain-Specific Results}

Table~\ref{tab:domain-results} breaks down performance by application domain.

\begin{table}[H]
\centering
\caption{Domain-Specific Results}
\label{tab:domain-results}
\begin{tabular}{lrrrrr}
\toprule
\textbf{Domain} & \textbf{Count} & \textbf{Verified} &
\textbf{Avg Coords} & \textbf{Avg Props} & \textbf{Avg Time} \\
\midrule
Data Structures & \expDomainDsCount & \expDomainDsVerified &
  \expDomainDsAvgCoords & \expDomainDsAvgProps & \suppDomainDsAvgTime \\
Mathematics & \expDomainMathCount & \expDomainMathVerified &
  \expDomainMathAvgCoords & \expDomainMathAvgProps & \suppDomainMathAvgTime \\
Sorting & \expDomainSortCount & \expDomainSortVerified &
  \expDomainSortAvgCoords & \expDomainSortAvgProps & \suppDomainSortAvgTime \\
String Processing & \expDomainStrCount & \expDomainStrVerified &
  \expDomainStrAvgCoords & \expDomainStrAvgProps & \suppDomainStrAvgTime \\
Web/API & \expDomainWebCount & \expDomainWebVerified &
  \expDomainWebAvgCoords & \expDomainWebAvgProps & \suppDomainWebAvgTime \\
State Machines & \expDomainSmCount & \expDomainSmVerified &
  \expDomainSmAvgCoords & \expDomainSmAvgProps & \suppDomainSmAvgTime \\
\bottomrule
\end{tabular}
\end{table}

Data Structures and State Machines exhibit the highest coordinate counts
(\expDomainDsAvgCoords\ and \expDomainSmAvgCoords\ respectively), reflecting
their richer internal structure.  Sorting algorithms have the fewest
coordinates (\expDomainSortAvgCoords), as they are typically single-function
programs.  Timing varies modestly across domains: the fastest domain
(String Processing, \suppDomainStrAvgTime) is within 20\% of the slowest
(State Machines, \suppDomainSmAvgTime).

\subsection{Test Suite Metrics by Program Size}

Table~\ref{tab:size-metrics} shows how test generation scales with program
size.

\begin{table}[H]
\centering
\caption{Test Suite Metrics by Program Size}
\label{tab:size-metrics}
\begin{tabular}{lrrrrr}
\toprule
\textbf{Category} & \textbf{Count} & \textbf{Mean Coords} &
\textbf{Mean Props} & \textbf{Mean Time} & \textbf{Mean Lines} \\
\midrule
Tiny & \compTinyCount & \compTinyMeanCoords & \compTinyMeanProps &
  \compTinyMeanTime & \compTinyMeanLines \\
Small & \compSmallCount & \compSmallMeanCoords & \compSmallMeanProps &
  \compSmallMeanTime & \compSmallMeanLines \\
Medium & \compMediumCount & \compMediumMeanCoords & \compMediumMeanProps &
  \compMediumMeanTime & \compMediumMeanLines \\
Large & \compLargeCount & \compLargeMeanCoords & \compLargeMeanProps &
  \compLargeMeanTime & \compLargeMeanLines \\
\bottomrule
\end{tabular}
\end{table}

Notably, generation time does not increase monotonically with program size.
Medium programs (\compMediumMeanTime) are slightly faster than small ones
(\compSmallMeanTime), likely because medium programs have more regular
structure that the site builder exploits for efficient overlap extraction.
The number of pairwise overlaps grows quadratically with coordinates:
from $\binom{\compTinyMeanCoords}{2} = 3$ for tiny programs to
$\binom{\compLargeMeanCoords}{2} = 105$ for large programs.

\subsection{Descent Strategy Comparison}

Table~\ref{tab:descent-strategies} compares the four descent strategies
on the comprehensive benchmark.

\begin{table}[H]
\centering
\caption{Descent Strategy Comparison}
\label{tab:descent-strategies}
\begin{tabular}{lrrrr}
\toprule
\textbf{Strategy} & \textbf{Runs} & \textbf{Successes} &
\textbf{Rate} & \textbf{Mean Time} \\
\midrule
Eager & \compDescentEagerRuns & \compDescentEagerSuccesses &
  \compDescentEagerRate & \compDescentEagerMeanMs\,ms \\
Exhaustive & \compDescentExhaustiveRuns & \compDescentExhaustiveSuccesses &
  \compDescentExhaustiveRate & \compDescentExhaustiveMeanMs\,ms \\
Iterative & \compDescentIterativeRuns & \compDescentIterativeSuccesses &
  \compDescentIterativeRate & \compDescentIterativeMeanMs\,ms \\
Optimistic & \compDescentOptimisticRuns & \compDescentOptimisticSuccesses &
  \compDescentOptimisticRate & \compDescentOptimisticMeanMs\,ms \\
\bottomrule
\end{tabular}
\end{table}

All four strategies achieve \compDescentEagerRate\ success rate on
our benchmark, with \compDescentEagerRuns\ runs each.  The Eager strategy
is recommended for test generation, as it finds the first obstruction
quickly (if any exist) and avoids unnecessary computation.  The Exhaustive
strategy is preferred for mutation analysis, where all obstructions must
be enumerated.

\subsection{Subsystem Profiling}

Table~\ref{tab:subsystem} profiles the key subsystems of the test
generation pipeline across program sizes.

\begin{table}[H]
\centering
\caption{Subsystem Profiling: Time by Program Size}
\label{tab:subsystem}
\begin{tabular}{lrrr}
\toprule
\textbf{Subsystem} & \textbf{Small} & \textbf{Medium} & \textbf{Large} \\
\midrule
Full Descent & \subRunfulldescentSmallTime &
  \subRunfulldescentMediumTime & \subRunfulldescentLargeTime \\
SMT Encoding & \subEncodeforsolverSmallTime &
  \subEncodeforsolverMediumTime & \subEncodeforsolverLargeTime \\
Spec Satisfaction & \subSpecificationsatisfactionSmallTime &
  \subSpecificationsatisfactionMediumTime &
  \subSpecificationsatisfactionLargeTime \\
Formal Core Site & \subFormalcoresiteSmallTime &
  \subFormalcoresiteMediumTime & \subFormalcoresiteLargeTime \\
Verification Orch. & \subOrchestrateverificationSmallTime &
  \subOrchestrateverificationMediumTime &
  \subOrchestrateverificationLargeTime \\
Theorem Economics & \subTheoremeconomicsSmallTime &
  \subTheoremeconomicsMediumTime & \subTheoremeconomicsLargeTime \\
\bottomrule
\end{tabular}
\end{table}

SMT encoding dominates for small programs (\subEncodeforsolverSmallTime)
due to solver initialization overhead, but amortizes well for medium and
large programs (\subEncodeforsolverMediumTime\ and
\subEncodeforsolverLargeTime\ respectively).  The full descent subsystem
maintains constant time across all sizes (\subRunfulldescentSmallTime\ to
\subRunfulldescentLargeTime), confirming that the descent check itself is
efficient relative to the site construction.

\subsection{Scalability Analysis}

We measure how the pipeline scales with the number of coordinates using
synthetic programs of increasing size.

\begin{table}[H]
\centering
\caption{Scalability: Time vs.\ Number of Coordinates}
\label{tab:scalability}
\begin{tabular}{rr}
\toprule
\textbf{Coordinates} & \textbf{Time} \\
\midrule
2 & \subScaleTwoTime \\
4 & \subScaleFourTime \\
8 & \subScaleEightTime \\
12 & \subScaleTwelveTime \\
16 & \subScaleSixteenTime \\
20 & \subScaleTwentyTime \\
\bottomrule
\end{tabular}
\end{table}

The scaling data shows that the pipeline handles up to 20 coordinates
with sub-millisecond descent times.  The theoretical quadratic bound
from Theorem~\ref{thm:complexity} is not observable in practice because
the constant factors are small and the SMT solver handles overlap queries
efficiently through incremental solving.

\subsection{Proposition and Coordinate Breakdown}

The supplementary analysis reveals the composition of propositions and
coordinates across the benchmark:

\begin{itemize}[noitemsep]
  \item \textbf{Propositions} (\suppPropTotal\ total):
    \suppPropStructuralCount\ structural (\suppPropStructuralPct),
    \suppPropBehavioralCount\ behavioral (\suppPropBehavioralPct),
    \suppPropRelationalCount\ relational (\suppPropRelationalPct),
    \suppPropResourceCount\ resource (\suppPropResourcePct).
  \item \textbf{Coordinates} (\suppCoordTotal\ total):
    \suppCoordFunctionCount\ function (\suppCoordFunctionPct),
    \suppCoordModuleCount\ module (\suppCoordModulePct),
    \suppCoordInterfaceCount\ interface (\suppCoordInterfacePct).
\end{itemize}

Function coordinates dominate (\suppCoordFunctionPct), and behavioral
propositions are most common (\suppPropBehavioralPct).  This distribution
aligns with the predominance of function-level testing in practice and
suggests that cover tests primarily validate behavioral consistency across
function boundaries.

\subsection{Threats to Validity}

\paragraph{Internal validity.}
Our benchmark programs are relatively small (up to $\compLargeMeanLines$
mean lines for the large category).  Larger industrial codebases may
exhibit different overlap structures and more complex proposition
languages that challenge the SMT-based input generation.

\paragraph{External validity.}
The \expDomainCount\ domains (Data Structures, Mathematics, Sorting,
String Processing, Web/API, State Machines) are representative of common
programming tasks but do not cover all application domains.  In particular,
concurrent and distributed programs---where descent obstructions are
expected to be more prevalent---are not represented.

\paragraph{Construct validity.}
Our sheaf-theoretic mutation score (Definition~\ref{def:sheaf-mutation})
differs from the classical mutation score.  While it captures
semantically meaningful mutations through the presheaf perturbation
framework, it may miss purely syntactic mutations that do not affect
the site structure.

\paragraph{Conclusion validity.}
All \compTotalPrograms\ programs in the comprehensive benchmark and all
\expTotalPrograms\ in the experiment benchmark achieved $\compOverallAccuracy$
accuracy.  The \compBuggyCount\ injected-bug programs were all detected
(\compBuggyFullPass\ full passes), providing confidence in the
bug-detection capability, though the injected faults are necessarily
simpler than real-world bugs.

% =====================================================================
\section{Related Work}\label{sec:related}
% =====================================================================

\paragraph{Coverage-based testing.}
Traditional coverage metrics~\cite{Zhu1997} measure syntactic exercise.
Our approach provides \emph{semantic} coverage through covering families.
Statement, branch, and path coverage~\cite{Ammann2016} are purely structural
and cannot detect integration failures where module interfaces disagree---the
key insight that motivates our sheaf-theoretic approach.

\paragraph{Property-based testing.}
QuickCheck~\cite{Claessen2000} and Hypothesis~\cite{MacIver2019} generate
random inputs satisfying properties.  Our cover tests derive properties
from sheaf-theoretic structure rather than user-written specifications.
The key advantage is that the covering family provides a \emph{complete}
set of properties (the restriction consistency conditions) without
requiring manual specification.

\paragraph{Concolic testing.}
Concolic (concrete-symbolic) testing~\cite{Sen2005,Godefroid2005} combines
concrete execution with symbolic reasoning to explore execution paths
systematically.  Our SMT-guided input generation
(Section~\ref{sec:smt-background}) shares the use of constraint solving
but operates on the topological structure of the semantic site rather than
on execution paths.  Concolic testing explores the control-flow graph;
our approach explores the covering family structure.

\paragraph{Fuzzing.}
Coverage-guided fuzzing~\cite{Zalewski2014,Boehme2017} generates test
inputs by mutating seeds and selecting for increased code coverage.
Our boundary-value analysis (Section~\ref{sec:boundary-analysis})
serves a similar role but is guided by the mathematical structure of
propositions rather than coverage feedback.  The two approaches are
complementary: fuzzing can fill the random-input budget in our pipeline.

\paragraph{Metamorphic testing.}
Metamorphic testing~\cite{Chen2018} checks \emph{metamorphic relations}
between outputs on related inputs.  The restriction consistency condition
$\res_{U_{ij}}^{U_i}(\mathcal{F}(U_i))(x) =
\res_{U_{ij}}^{U_j}(\mathcal{F}(U_j))(x)$ is a metamorphic relation
in this sense: it relates the outputs of two coordinates on the same
input.  Our approach systematically derives these relations from the
site structure rather than requiring manual identification.

\paragraph{Specification-based testing.}
Formal specification-based testing~\cite{Hierons2009} generates tests from
formal specifications (pre/postconditions, invariants).  \JuGeo's judgment
presheaf $\mathcal{F}$ can be viewed as a specification in this framework,
with propositions serving as the formal contracts.  Our approach
automatically constructs the ``specification'' from program analysis,
avoiding the overhead of manual specification writing.

\paragraph{Formal test generation.}
Symbolic execution~\cite{King1976,Cadar2008} generates tests from path
constraints.  Our approach uses the topological structure of the semantic
site rather than execution paths, and mutation testing~\cite{Jia2011}
complements obstruction-derived regression tests.  The key distinction
is that symbolic execution targets path coverage, while our approach
targets \emph{semantic coverage} of the covering family---a fundamentally
different adequacy criterion.

% =====================================================================
\section{Conclusion}\label{sec:conclusion}
% =====================================================================

We have presented a sheaf-theoretic approach to test suite generation
exploiting covering families and descent obstructions.  The framework
rests on several key theoretical results:

\begin{itemize}[noitemsep]
  \item \textbf{Cover Test Completeness} (Theorem~\ref{thm:completeness}):
    when $\Hcoh{1} = 0$, passing all cover tests implies descent success.
  \item \textbf{Quantitative Completeness}
    (Theorem~\ref{thm:quant-completeness}): at most
    $\binom{n}{2} \cdot (2dM+1)^d$ tests suffice.
  \item \textbf{Finite Witness Theorem} (Theorem~\ref{thm:finite-witness}):
    finitely many inputs detect all obstruction classes.
  \item \textbf{Obstruction Witness Existence}
    (Theorem~\ref{thm:obstr-witness}): non-trivial $\Hcoh{1}$ always
    yields concrete witnesses.
  \item \textbf{Mutation Detection Completeness}
    (Theorem~\ref{thm:mutation-complete}): cover tests kill all
    boundary mutations.
  \item \textbf{Minimality} (Corollary~\ref{cor:minimality}): the
    obstruction suite has no redundant tests.
\end{itemize}

Evaluation on \expTotalPrograms\ programs across \expDomainCount\ domains
demonstrates practical effectiveness, with the comprehensive suite of
\compTotalPrograms\ programs achieving \compOverallAccuracy\ structural
accuracy.  Tests detected all \compBuggyCount\ injected-fault programs
with mean proposition verification ratio \compBuggyMeanPropRatio,
confirming strong bug detection capability.

\paragraph{Future directions.}
Several extensions merit investigation.  First, \emph{higher-order cover
tests} from $\Hcoh{n}$ for $n \geq 2$ would target multi-component
integration failures beyond pairwise interactions
(Section~\ref{sec:higher-obstr}).  Second, integration with \emph{CI/CD
pipelines} would enable continuous sheaf-theoretic testing, using the
incremental update algorithm (Algorithm~\ref{alg:incremental}) to
minimize regeneration cost on each commit.  Third, extending the
framework to \emph{concurrent programs} where descent obstructions
capture race conditions and synchronization failures would address a
major gap in current testing practice.  Finally, combining our
boundary-value analysis with coverage-guided fuzzing could yield a
hybrid approach that inherits the mathematical guarantees of our
framework with the practical scalability of fuzzing.

\paragraph{Acknowledgments.}
We thank the Z3 team for the SMT solver infrastructure, and the Lean
community for the proof assistant ecosystem.

% ─────────────────────────────────────────────────────────────────────
\begin{thebibliography}{99}

\bibitem{Zhu1997}
H.~Zhu, P.~Hall, and J.~May, ``Software Unit Test Coverage and Adequacy,''
ACM Computing Surveys, vol.~29, no.~4, 1997.

\bibitem{Claessen2000}
K.~Claessen and J.~Hughes, ``QuickCheck: A Lightweight Tool for Random
Testing of Haskell Programs,'' ICFP 2000.

\bibitem{MacIver2019}
D.~MacIver et al., ``Hypothesis: A New Approach to Property-Based Testing,''
Journal of Open Source Software, 2019.

\bibitem{Jia2011}
Y.~Jia and M.~Harman, ``An Analysis and Survey of the Development of
Mutation Testing,'' IEEE TSE, vol.~37, no.~5, 2011.

\bibitem{King1976}
J.~King, ``Symbolic Execution and Program Testing,'' CACM, vol.~19, no.~7,
1976.

\bibitem{Cadar2008}
C.~Cadar, D.~Dunbar, and D.~Engler, ``KLEE: Unassisted and Automatic
Generation of High-Coverage Tests,'' OSDI 2008.

\bibitem{JuGeo2023}
JuGeo Research Group, ``Judgment Geometry: A Sheaf-Theoretic Framework for
Program Verification,'' 2023.

\bibitem{Lean4}
L.~de~Moura et al., ``The Lean 4 Theorem Prover and Programming Language,''
CADE 2021.

\bibitem{Z32008}
L.~de~Moura and N.~Bjørner, ``Z3: An Efficient SMT Solver,'' TACAS 2008.

\bibitem{Ammann2016}
P.~Ammann and J.~Offutt, \emph{Introduction to Software Testing},
Cambridge University Press, 2nd edition, 2016.

\bibitem{Sen2005}
K.~Sen, D.~Marinov, and G.~Agha, ``CUTE: A Concolic Unit Testing Engine
for C,'' ESEC/FSE 2005.

\bibitem{Godefroid2005}
P.~Godefroid, N.~Klarlund, and K.~Sen, ``DART: Directed Automated Random
Testing,'' PLDI 2005.

\bibitem{Zalewski2014}
M.~Zalewski, ``American Fuzzy Lop (AFL),'' Technical Report, 2014.

\bibitem{Boehme2017}
M.~B\"ohme, V.-T.~Pham, and A.~Roychoudhury, ``Coverage-Based Greybox
Fuzzing as Markov Chain,'' IEEE TSE, vol.~45, no.~5, 2017.

\bibitem{Chen2018}
T.~Y.~Chen, F.-C.~Kuo, H.~Liu, P.-L.~Poon, D.~Towey, T.~H.~Tse, and
Z.~Q.~Zhou, ``Metamorphic Testing: A Review of Challenges and
Opportunities,'' ACM Computing Surveys, vol.~51, no.~1, 2018.

\bibitem{Hierons2009}
R.~M.~Hierons et al., ``Using Formal Specifications to Support Testing,''
ACM Computing Surveys, vol.~41, no.~2, 2009.

\end{thebibliography}

\end{document}
